\chapter{The Probabilistic Method}

\begin{center}\itshape
To prove that a haystack contains a needle,\\
show that a straw picked at random is a needle with positive probability.
\end{center}

Some objects are easy to describe and hard to find: a colouring of a large complete graph with no big monochromatic
clique, an assignment satisfying a formula whose clauses barely overlap, a sparse graph with a giant connected
piece. Paul Erd\H os discovered in 1947 that one can prove such objects exist without constructing them: choose the
object at random, and show that it has the required property with positive probability. Since then the
\emph{probabilistic method} has become one of the main tools of combinatorics and theoretical computer science.

The chapter climbs in sophistication. The \emph{first moment method} uses only linearity of expectation: if the expected
number of bad configurations is less than $1$, some outcome has none. \emph{Alterations} first take a random object, then
repair it. The \emph{second moment method} shows that a random variable is positive with high probability, which
locates thresholds in random graphs. The \emph{Lov\'asz local lemma} handles many bad events that are each rare and
depend on few others, even when the union bound fails. Then we follow the random graph $G(n,\frac cn)$ through its phase
transition: for $c<1$ all components are small, for $c>1$ a giant one appears, and the branching processes of Chapter~3
explain why. Finally, \emph{conditional expectations} turn some of these existence proofs into fast deterministic
algorithms.

\begin{think}
Among any $6$ people there are $3$ mutual friends or $3$ mutual strangers. Is there a party of $1000$ people with no $20$
mutual friends and no $20$ mutual strangers? If so, could you describe the friendships explicitly?
\end{think}

\section{The first moment method}

Two facts of Chapter~2 do all the work: if $X$ is a random variable with values in $\N$ and $\mathbb E[X]<1$, then $P(X=0)>0$; and any random
variable takes, with positive probability, a value at least its mean. So an object with the desired property exists.

The \emph{Ramsey number} $R(k)$ is the smallest $n$ such that every colouring of the edges of the complete graph $K_n$ in red and blue
contains a monochromatic $K_k$ (a set of $k$ vertices whose $\binom k2$ edges all have the same colour).

\begin{proposition}\label{prop:ramseyupper}
$R(3)=6$, and for all $k\ge2$, $R(k)\le\binom{2k-2}{k-1}<4^{k-1}$.
\end{proposition}
\begin{proof}
Let $R(k,l)$ be the smallest $n$ such that every red--blue colouring of $K_n$ has a red $K_k$ or a blue $K_l$. We show $R(k,l)\le R(k-1,l)+R(k,l-1)$; with
$R(k,2)=k$ and $R(2,l)=l$, induction and Pascal's rule give $R(k,l)\le\binom{k+l-2}{k-1}$. Take $n=R(k-1,l)+R(k,l-1)$ and a vertex $v$ of $K_n$. Among the other
$n-1$ vertices, either at least $R(k-1,l)$ are joined to $v$ in red, or at least $R(k,l-1)$ in blue. In the first case, these vertices contain a blue $K_l$, or
a red $K_{k-1}$ which together with $v$ forms a red $K_k$; the second case is symmetric. For $k=3$ this gives $R(3)\le6$, and the colouring of $K_5$
of Figure~\ref{fig:ramsey} (the sides of a pentagon red, the diagonals blue) has no monochromatic triangle, so $R(3)=6$.
\end{proof}

\begin{figure}[ht]
\centering
\fig{ch12-ramsey}
\caption{Left: a colouring of $K_5$ without a monochromatic triangle. Right: a random colouring of $K_6$; it has a monochromatic
triangle (thick), as every colouring of $K_6$ must.}
\label{fig:ramsey}
\end{figure}

\begin{theorem}[Erd\H os, 1947]\label{thm:erdosramsey}
If $\binom nk2^{1-\binom k2}<1$, then $R(k)>n$. In particular $R(k)>2^{k/2}$ for $k\ge3$.
\end{theorem}
\begin{proof}
Colour each edge of $K_n$ red or blue independently with probability $\frac12$. A given set of $k$ vertices is monochromatic with probability
$2\cdot2^{-\binom k2}$, so the expected number of monochromatic $K_k$ is $\binom nk2^{1-\binom k2}<1$. Hence some colouring has none, and $R(k)>n$. For $n=\lfloor2^{k/2}\rfloor$,
\[
\begin{aligned}\binom nk2^{1-\binom k2}&\le\frac{n^k}{k!}2^{1-\frac{k(k-1)}2}\eqbr\le\frac{2^{k^2/2}}{k!}2^{1-\frac{k^2}2+\frac k2}=\frac{2^{1+k/2}}{k!},\end{aligned}
\]
which is $\frac{2^{2.5}}6<1$ for $k=3$ and decreases with $k$, since each step multiplies it by $\frac{\sqrt2}{k+1}<1$.
\end{proof}

So $\sqrt2^{\,k}<R(k)<4^k$. For $k=10$ the theorem gives $R(10)>100$, and for $k=20$, $R(20)>5817$. Nobody knows how to write down an explicit colouring
anywhere near these bounds: the best explicit constructions are far smaller for large $k$, and Erd\H os offered a prize for one of
size $(1+\varepsilon)^k$. Exact values are known only for $k\le4$ ($R(4)=18$); for $k=5$ it is known that $43\le R(5)\le46$. The upper bound $4^k$ was improved to
$(4-\varepsilon)^k$ only in 2023, by Campos, Griffiths, Morris and Sahasrabudhe. Volume~3, Appendix~B takes up Ramsey numbers from the combinatorial side:
the proof that $R(4)=18$, Tur\'an's theorem on dense graphs without cliques, and a probabilistic proof of it in the spirit of this chapter.

\begin{proposition}[large cuts]\label{prop:maxcut}
Every graph with $m$ edges has a \emph{cut} with at least $\frac m2$ edges: its vertices can be split into two sets so that at least $\frac m2$
edges join the two sets.
\end{proposition}
\begin{proof}
Put each vertex in one of the two sets by an independent fair coin. Each edge is cut with probability $\frac12$, so the expected
number of cut edges is $\frac m2$, and some split achieves at least this.
\end{proof}

\section{Alterations}

A random object may have a few defects; if they are few enough, we remove them.

\begin{proposition}\label{prop:alteration}
\begin{enumerate}[label=(\alph*)]
\item For every $n$, $R(k)>n-\binom nk2^{1-\binom k2}$.
\item A graph with $n$ vertices and average degree $d\ge1$ has an independent set (a set of vertices with no edge inside) of size at
least $\frac n{2d}$.
\end{enumerate}
\end{proposition}
\begin{proof}
(a) Colour $K_n$ at random and delete one vertex from each monochromatic $K_k$. What remains has no monochromatic $K_k$, and its expected
number of vertices is at least $n-\binom nk2^{1-\binom k2}$; some outcome does at least as well, and the number of remaining vertices is an integer.

(b) Keep each vertex with probability $p$, independently, and then delete one endpoint of every edge with both ends kept. The rest is
independent. The graph has $\frac{nd}2$ edges, so the expected number of vertices left is at least $pn-p^2\frac{nd}2$; with $p=\frac1d$ this is $\frac n{2d}$.
\end{proof}

For $k=10$, part (a) with $n=127$ gives $R(10)>115$, against $100$ from Theorem~\ref{thm:erdosramsey}; asymptotically, the alteration gains a factor $\sqrt2$.
Part (b) is a weak form of Tur\'an's theorem; Exercise~\ref{ex:carowei} gives a sharper one with a cleverer random choice.

\section{The second moment method}

The first moment shows that $X=0$ is possible. To show that $X>0$ is likely, we need the variance.

\begin{lemma}[second moment]\label{lem:secondmoment}
If $X$ is a random variable with $\mathbb E[X^2]<\infty$ and $\mathbb E[X]>0$, then $P(X=0)\le\frac{\operatorname{Var}(X)}{\mathbb E[X]^2}$.
\end{lemma}
\begin{proof}
If $X=0$ then $|X-\mathbb E X|\ge\mathbb E X$; apply Chebyshev's inequality.
\end{proof}

The \emph{random graph} $G(n,p)$ has vertex set $\{1,\dots,n\}$, and each of the $\binom n2$ possible edges is present independently with probability $p$.
As $p$ grows from $0$ to $1$, the graph acquires properties one after another, and many of them appear suddenly. We
say that an event holds \emph{with high probability} (w.h.p.) if its probability tends to $1$ as $n\to\infty$; here $p$ may depend on $n$.

\begin{theorem}[threshold for triangles]\label{thm:triangles}
Let $T$ be the number of triangles in $G(n,p)$. If $np\to0$, then w.h.p.\ $T=0$. If $np\to\infty$, then w.h.p.\ $T>0$.
\end{theorem}
\begin{proof}
For each set $A$ of three vertices let $I_A$ be the indicator that its three edges are present, so $T=\sum_AI_A$ and $\mathbb E[T]=\binom n3p^3\le\frac{(np)^3}6$. If $np\to0$,
then $P(T>0)\le\mathbb E[T]\to0$ by Markov's inequality.

Now let $np\to\infty$, so $\mathbb E[T]\sim\frac{(np)^3}6\to\infty$. Expand $\operatorname{Var}(T)=\sum_{A,B}\operatorname{Cov}(I_A,I_B)$. If $A$ and $B$ share at most one vertex, they have no edge in
common, so $I_A$ and $I_B$ are independent and the covariance is $0$. If $A=B$, it is at most $\mathbb E[I_A]=p^3$. If $A\ne B$ share two vertices, then
$A\cup B$ has four vertices and five edges, and the covariance is at most $\mathbb E[I_AI_B]=p^5$. Each set of four vertices contains four triangles, hence
$12$ ordered pairs of distinct ones, all sharing an edge. So
\[
\operatorname{Var}(T)\le\mathbb E[T]+12\binom n4p^5,
\]
and by the lemma
\[
\begin{aligned}P(T=0)&\le\frac1{\mathbb E[T]}+\frac{12\binom n4p^5}{\binom n3^2p^6}\eqbr\le\frac1{\mathbb E[T]}+\frac{18+o(1)}{n^2p}\to0,\end{aligned}
\]
since $\binom n4\big/\binom n3^2\sim\frac{36}{24n^2}$ and $n^2p\ge np\to\infty$.
\end{proof}

The function $\frac1n$ is a \emph{threshold} for containing a triangle. At the threshold itself, when $np=c$, the number of triangles converges in law to a
Poisson variable of mean $\frac{c^3}6$ (by the method of factorial moments of Chapter~2; see Alon and Spencer, Chapter~8), so $P(T>0)\to1-e^{-c^3/6}$. Every property preserved by adding edges
has a threshold (Bollob\'as and Thomason, 1987), and Exercise~\ref{ex:connectivity} finds the one for connectivity, $\frac{\ln n}n$.

\section{The local lemma}

If $A_1,\dots,A_n$ are bad events, the union bound shows that none occurs with positive probability when $\sum P(A_i)<1$. If they are
independent, it suffices that each $P(A_i)<1$. The local lemma interpolates: it suffices that each event is rare compared with the
number of events it depends on.

An event $A$ is \emph{mutually independent} of a family of events $\{A_j:j\in J\}$ if $P(A\cap C)=P(A)P(C)$ for every $C$ that is an intersection of events
$A_j$ or complements $A_j^c$, $j\in J$. Write $\bar A=A^c$.

\begin{theorem}[Lov\'asz local lemma, symmetric form]\label{thm:lll}
Let $A_1,\dots,A_n$ be events with $P(A_i)\le p$ for all $i$. Suppose that each $A_i$ is mutually independent of all but at most $d$ of the other
events, and that $ep(d+1)\le1$. Then
\[
P\Bigl(\bigcap_{i=1}^n\bar A_i\Bigr)\ge\Bigl(1-\frac1{d+1}\Bigr)^n>0.
\]
\end{theorem}
\begin{proof}
For each $i$, let $D_i$ be a set of at most $d$ indices such that $A_i$ is mutually independent of $\{A_j:j\notin D_i\cup\{i\}\}$. For $S\subseteq\{1,\dots,n\}$ write
$B_S=\bigcap_{j\in S}\bar A_j$ (with $B_\emptyset=\Omega$). We prove by induction on $|S|$ that $P(B_S)>0$ and that
\[
P(A_i\mid B_S)\le\frac1{d+1}\qquad\text{for every }i\notin S.\tag{$*$}
\]
For $S=\emptyset$, $P(A_i)\le p\le\frac1{e(d+1)}$. Let $|S|\ge1$ and suppose the claims hold for smaller sets. First, removing the elements of $S$ one at a
time, $P(B_S)=\prod_tP(\bar A_{s_t}\mid B_{\{s_1,\dots,s_{t-1}\}})\ge(1-\frac1{d+1})^{|S|}>0$ by $(*)$ for the smaller sets. Now fix $i\notin S$ and split $S$ into $S_1=S\cap D_i$ and $S_2=S\setminus D_i$.
If $S_1=\emptyset$, then $P(A_i\mid B_S)=P(A_i)\le p$ by mutual independence. Otherwise
\[
P(A_i\mid B_S)=\frac{P(A_i\cap B_{S_1}\mid B_{S_2})}{P(B_{S_1}\mid B_{S_2})}.
\]
The numerator is at most $P(A_i\mid B_{S_2})=P(A_i)\le p$, by mutual independence. For the denominator, list $S_1=\{j_1,\dots,j_r\}$ with $r\le d$; by $(*)$ applied to the
sets $S_2\cup\{j_1,\dots,j_{t-1}\}$, which are smaller than $S$,
\[
\begin{aligned}P(B_{S_1}\mid B_{S_2})&=\prod_{t=1}^r\Bigl(1-P\bigl(A_{j_t}\mid B_{S_2\cup\{j_1,\dots,j_{t-1}\}}\bigr)\Bigr)\eqbr\ge\Bigl(1-\frac1{d+1}\Bigr)^d\ge\frac1e,\end{aligned}
\]
since $(1+\frac1d)^d\le e$. So $P(A_i\mid B_S)\le ep\le\frac1{d+1}$. Finally, $P(B_{\{1,\dots,n\}})\ge(1-\frac1{d+1})^n$ as above.
\end{proof}

A \emph{$k$-CNF formula} is a conjunction of clauses, each a disjunction of $k$ literals (variables or their negations) on $k$ distinct
variables, such as $(x_1\vee\bar x_2\vee x_4)\wedge(\bar x_1\vee x_3\vee\bar x_4)$. Deciding whether such a formula can be satisfied is the archetypal NP-complete problem for
$k\ge3$.

\begin{corollary}\label{cor:ksat}
If each clause of a $k$-CNF formula shares a variable with at most $d$ other clauses, and $e(d+1)\le2^k$, then the formula is satisfiable.
\end{corollary}
\begin{proof}
Give the variables independent uniform values in $\{\text{true},\text{false}\}$, and let $A_i$ be the event that clause $i$ is false: all its $k$ literals are false,
probability $2^{-k}$. The event $A_i$ is determined by the variables of clause $i$, so it is mutually independent of the events of the clauses
that share no variable with it (Proposition~\ref{prop:indepfun}, functions of disjoint groups of independent variables). Theorem~\ref{thm:lll} with $p=2^{-k}$ applies.
\end{proof}

The union bound would require fewer than $2^k$ clauses in total; the local lemma allows any number, as long as each clause
meets fewer than $\frac{2^k}e$ others. For a long time the lemma gave no way to find the satisfying assignment, which is exponentially rare.
In 2009 Robin Moser found an astonishingly simple algorithm: start from a random assignment, and as long as some clause
is false, resample its variables at random. Moser and Tardos (2010) proved that under the hypothesis of the local
lemma the expected number of resamplings is at most $\frac n d$, where $n$ is the number of clauses. A numerical calculation runs it.

\section{The phase transition of random graphs}

Now let $p=\frac cn$ with a constant $c>0$: each vertex has about $c$ neighbours on average. Erd\H os and R\'enyi discovered in 1960 that
the structure of $G(n,\frac cn)$ changes abruptly at $c=1$ (Figure~\ref{fig:giant}). The key is to explore the component $C(v)$ of a vertex $v$ step by step.

\emph{Exploration.} Vertices are \emph{active}, \emph{explored} or \emph{neutral}. Initially $v$ is active and all others are neutral. At each step $i=1,2,\ldots$,
choose an active vertex $w$, reveal its edges to the neutral vertices, make the $X_i$ neutral neighbours found active, and mark
$w$ explored. The number of active vertices after step $i$ is $A_i=1+X_1+\dots+X_i-i$, the exploration stops at the first $i$ with $A_i=0$, and then $|C(v)|=i$.
Each edge is revealed at most once, so given the past, $X_i$ is binomial$(N_i,p)$, where $N_i\le n-1$ is the number of neutral vertices. Just as for a
Galton--Watson process with Poisson$(c)$ offspring (Chapter~3), the exploration should die out if $c<1$ and may survive if $c>1$.

\begin{theorem}[small components]\label{thm:subcritical}
Let $c<1$ and $I=c-1-\ln c>0$. For every $a>\frac1I$, w.h.p.\ every component of $G(n,\frac cn)$ has at most $a\ln n$ vertices.
\end{theorem}
\begin{proof}
\emph{Coupling.} At step $i$, add to the $N_i$ trials of the exploration $n-N_i$ independent ``ghost'' trials of probability $p$, and let $\tilde X_i\ge X_i$ count the
successes among all $n$ trials (after the exploration has stopped, let $\tilde X_i$ count $n$ fresh trials). Given the past, $\tilde X_i$ is binomial$(n,p)$, so the $\tilde X_i$
are independent binomial$(n,p)$ variables, and $\tilde X_1+\dots+\tilde X_k$ is binomial$(nk,p)$, with mean $ck$.

\emph{One vertex.} If $|C(v)|>k$, the exploration is still running after $k$ steps, so $A_k\ge1$, that is $X_1+\dots+X_k\ge k$, and so $\tilde X_1+\dots+\tilde X_k\ge k$. By the
Chernoff bound (Theorem~\ref{thm:chernoff}) with $\mu=ck$ and $1+\delta=\frac1c$,
\[
\begin{aligned}P(|C(v)|>k)&\le\Bigl(\frac{e^\delta}{(1+\delta)^{1+\delta}}\Bigr)^{ck}\eqbr=e^{-k(c-1-\ln c)}=e^{-kI}.\end{aligned}
\]
\emph{All vertices.} With $k=\lfloor a\ln n\rfloor$, the union bound over the $n$ vertices gives
\[
P\Bigl(\max_v|C(v)|>k\Bigr)\le ne^{-kI}\le e^In^{1-aI}\to0.
\]
\end{proof}

For $c>1$ the exploration from a vertex survives with probability close to the survival probability $\rho$ of the Poisson$(c)$ branching
process of Exercise~\ref{ex:poissongw}, and the surviving explorations all end up in a single component.

\begin{remark}
For $c>1$, w.h.p.\ $G(n,\frac cn)$ has a unique \emph{giant component}, with $(\rho+o(1))n$ vertices, where $\rho\in\,]0,1[$ solves $1-\rho=e^{-c\rho}$, while all other
components have $O(\ln n)$ vertices (Erd\H os and R\'enyi, 1960). Exercise~\ref{ex:giant} proves the existence of a component of linear size by a
depth-first search argument of Krivelevich and Sudakov; the uniqueness and the exact constant $\rho$ need finer couplings, found in
the books of Janson, \L uczak and Ruci\'nski, of Frieze and Karo\'nski, and of van der Hofstad. A numerical calculation watches $\rho$ appear. At $c=1$
itself the largest component has about $n^{2/3}$ vertices, a critical window studied by Aldous with Brownian motion.
\end{remark}

\begin{figure}[ht]
\centering
\fig{ch12-giant}
\caption{The largest (blue) and second largest (orange) components of a random graph on $n=20\,000$ vertices with $\frac{cn}2$ random edges (essentially
$G(n,\frac cn)$), as fractions of $n$, against $c$; edges are added one by one, so that all values of $c$ come from one growing graph.
The black curve is the survival probability $\rho(c)$ of a Poisson$(c)$ branching process.}
\label{fig:giant}
\end{figure}

\begin{beauty}
A random graph knows about branching processes. Seen from one vertex, $G(n,\frac cn)$ looks like a Poisson$(c)$ family tree: the number
of neighbours is close to Poisson$(c)$, and so is the number of new neighbours of each of them, as long as the exploration is
small compared with $n$. The extinction probability of Galton's family names (Chapter~3) becomes the fraction of
vertices outside the giant component.
\end{beauty}

\section{From existence to algorithms}

A probabilistic proof often hides an algorithm. Choose the random bits one at a time, each time keeping the
conditional expectation of the quantity of interest from decreasing. When these conditional expectations can be
computed, this is the \emph{method of conditional expectations}.

\begin{theorem}\label{thm:maxsat}
For a $k$-CNF formula with $m$ clauses, there is an assignment satisfying at least $(1-2^{-k})m$ clauses, and the following deterministic
algorithm finds one: for $i=1,\dots,n$, give $x_i$ the value that maximises the expected number of satisfied clauses when $x_1,\dots,x_i$ are fixed
and the remaining variables are uniform.
\end{theorem}
\begin{proof}
Let $W$ be the number of satisfied clauses under a uniform random assignment and $W_i(a_1,\dots,a_i)$ its conditional expectation given
$x_1=a_1,\dots,x_i=a_i$, that is, $\mathbb E[W\mid x_1,\dots,x_i]$ evaluated at these values (Chapter~10, for a finite partition). A clause contributes $1$ if one of its fixed
literals is true, and otherwise $1-2^{-u}$, where $u$ is its number of unfixed literals ($0$ if $u=0$): each $W_i$ is computable in time proportional
to the size of the formula. By the tower property, writing $a=(a_1,\dots,a_{i-1})$,
\[
\begin{aligned}W_{i-1}(a)&=\tfrac12W_i(a,\text{true})\eqbr+\tfrac12W_i(a,\text{false}),\end{aligned}
\]
so the better of the two values satisfies $W_i\ge W_{i-1}$. Starting from $W_0=\mathbb E[W]=(1-2^{-k})m$, the algorithm ends with a full assignment satisfying
$W_n\ge(1-2^{-k})m$ clauses.
\end{proof}

For $k=3$ this gives $\frac78$ of the clauses. H\aa stad proved in 2001 that, unless P $=$ NP, no polynomial algorithm can guarantee $\frac78+\varepsilon$ for
satisfiable $3$-CNF formulas: the trivial random assignment, derandomised, is optimal. For MAX-CUT, the same method
places the vertices one by one on the side that cuts more edges to the already placed ones, and finds a cut of at least
$\frac m2$ edges (Proposition~\ref{prop:maxcut}).

\section{Exercises}

\begin{exercise}[$\star$]\label{ex:carowei}
\emph{The Caro--Wei bound.} Order the vertices of a graph uniformly at random, and keep the vertices that come before all their
neighbours. Show that this is an independent set, and deduce that the graph has an independent set of size at least
$\sum_v\frac1{d(v)+1}$, where $d(v)$ is the degree of $v$. Compare with Proposition~\ref{prop:alteration}(b).
\end{exercise}

\begin{exercise}[$\star$]\label{ex:tournament}
\emph{Tournaments.} In a round-robin tournament of $n$ players every pair plays once, with no draws. Show that if
$\binom nk(1-2^{-k})^{n-k}<1$, there is a tournament in which every set of $k$ players is beaten by some other player (who beats all $k$ of
them). Deduce that such tournaments exist for $k=2$ with $n=21$ players and for $k=3$ with $n=91$.
\end{exercise}

\begin{exercise}[$\star\star$]\label{ex:sumfree}
\emph{Sum-free sets (Erd\H os, 1965).} A set of integers is \emph{sum-free} if it contains no $x$, $y$, $z$ with $x+y=z$. Show that every set $A$ of $n$
non-zero integers contains a sum-free subset of more than $\frac n3$ elements. (Take a prime $p=3q+2$ larger than all $|a|$, $a\in A$; for a uniform
$t\in\{1,\dots,p-1\}$, keep the $a\in A$ such that $ta\bmod p$ lies in the middle third $\{q+1,\dots,2q+1\}$.)
\end{exercise}

\begin{exercise}[$\star\star$]\label{ex:crossing}
\emph{The crossing lemma.} The crossing number $\operatorname{cr}(G)$ of a graph is the least number of edge crossings in a drawing of it in the plane.
Using Euler's formula for planar graphs (a planar graph with $n\ge3$ vertices has at most $3n-6$ edges), show that $\operatorname{cr}(G)\ge m-3n$ for
every graph with $n$ vertices and $m$ edges. Then apply this to a random induced subgraph to show that $\operatorname{cr}(G)\ge\frac{m^3}{64n^2}$ when $m\ge4n$.
\end{exercise}

\begin{exercise}[$\star$]\label{ex:propertyb}
\emph{Colouring hypergraphs.} A family of $k$-element sets is \emph{2-colourable} if the elements can be coloured red and blue so that no set
of the family is monochromatic. Show that a family of fewer than $2^{k-1}$ sets is 2-colourable, and that so is any family in which each
set meets at most $d$ others, if $e(d+1)\le2^{k-1}$.
\end{exercise}

\begin{exercise}[$\star\star\star$]\label{ex:connectivity}
\emph{The connectivity threshold.} Let $p=\frac{\ln n+\omega}n$, where $\omega=\omega(n)$.
(a) Let $Y$ be the number of isolated vertices. Show that $\mathbb E[Y]\sim e^{-\omega}$ if $\omega=o(\ln n)$, and that if $\omega\to-\infty$, then w.h.p.\ $Y>0$, so the graph is not connected.
(b) Suppose $\omega\to+\infty$ with $\omega\le\ln\ln n$. Show that the expected number of sets $S$ of $k$ vertices, $1\le k\le\frac n2$, with no edge between $S$ and its complement
and spanning a connected subgraph, tends to $0$. Conclude that w.h.p.\ the graph is connected. (A connected graph on $k$ vertices contains a
spanning tree, and there are $k^{k-2}$ trees on $k$ labelled vertices, by Cayley's formula.) Explain why the condition $\omega\le\ln\ln n$ may be dropped.
\end{exercise}

\begin{exercise}[$\star\star\star$]\label{ex:giant}
\emph{A giant component, by depth-first search.} Let $p=\frac{1+\varepsilon}n$ with $0<\varepsilon\le\frac1{10}$. Run a depth-first search on $G(n,p)$, maintaining three sets: $S$
(finished vertices), $U$ (a stack) and $T$ (unvisited vertices). If $U$ is empty, move some vertex of $T$ to $U$. Otherwise, let $v$ be the top of
$U$, and query the pairs $\{v,w\}$, $w\in T$, one at a time: at the first edge found, push $w$ onto $U$; if there is none, move $v$ to $S$. The answers to
the queries are independent Bernoulli$(p)$ variables.
(a) Show that the vertices of $U$ always lie on a path of the graph, that every pair with one vertex in $S$ and one in $T$ has been
queried, and that $|S|+|U|$ is at least the number of positive answers so far.
(b) Let $N=\lfloor\frac{\varepsilon n^2}2\rfloor$. Show that w.h.p.\ the first $N$ answers contain at least $\frac{\varepsilon(1+\varepsilon)n}2-n^{2/3}$ positive ones.
(c) Show that after $N$ queries, $|S|<\frac n3$. (Otherwise consider the moment when $|S|=\frac n3$.)
(d) Deduce that after $N$ queries, w.h.p.\ $|U|\ge\frac{\varepsilon^2n}5$. Conclude that w.h.p.\ $G(n,p)$ has a component, and even a path, with at least $\frac{\varepsilon^2n}5$ vertices,
and that the same holds for $G(n,\frac cn)$ for every $c>1$.
\end{exercise}

\begin{exercise}[$\star\star$]\label{pb:vanish}
\begin{enumerate}[label=(\alph*)]
\item Let $p$ be prime and $\zeta=e^{2i\pi/p}$. Show that rationals $a_0,\dots,a_{p-1}$ satisfy $\sum_ka_k\zeta^k=0$ if and
only if they are all equal. Which sets of $p$-th roots of unity have sum $0$?
\item Find all non-empty sets of sixth roots of unity with sum $0$. (There are $9$.)
\item Let $n=p^a$ be a prime power. Show that a set of $n$-th roots of unity has sum $0$ if and only if it is a union of
sets of the form $\{\eta,\eta\omega,\dots,\eta\omega^{p-1}\}$, with $\omega=e^{2i\pi/p}$.
\end{enumerate}
\end{exercise}

\begin{exercise}[$\star$]\label{pb:samecycle}
Let $\sigma$ be a uniformly random permutation of $\{1,\dots,n\}$.
\begin{enumerate}[label=(\alph*)]
\item Show that the length of the cycle containing $1$ is uniform on $\{1,\dots,n\}$.
\item Show that $1,2,\dots,k$ lie in the same cycle with probability $\frac1k$, for $1\le k\le n$.
\end{enumerate}
\end{exercise}

\begin{exercise}[$\star\star$]\label{pb:denes}
\begin{enumerate}[label=(\alph*)]
\item Let $\tau_1,\dots,\tau_{n-1}$ be transpositions of $\{1,\dots,n\}$, $\tau_i=(a_i\,b_i)$. Show that $\tau_1\tau_2\cdots\tau_{n-1}$ is an
$n$-cycle if and only if the edges $\{a_i,b_i\}$ form a tree.
\item (Dénes, 1959.) Deduce that $(1\,2\,\dots\,n)$ can be written as a product of $n-1$ transpositions in exactly $n^{n-2}$
ways.
\end{enumerate}
\end{exercise}

\begin{exercise}[$\star$]\label{pb:conics}
\emph{Conics.} Let $p$ be an odd prime.
\begin{enumerate}[label=(\alph*)]
\item Let $a\ne0$ and $b,c$ be in $\mathbb F_p$, and $D=b^2-4ac$. Show that
\[
\sum_{x\in\mathbb F_p}\chi(ax^2+bx+c)=\begin{cases}-\chi(a)&\text{if }D\ne0,\\(p-1)\chi(a)&\text{if }D=0.\end{cases}
\]
\item Let $Q(x,y)=ax^2+bxy+cy^2$ with $D=b^2-4ac\ne0$. Show that the equation $Q(x,y)=1$ has exactly $p-\chi(D)$ solutions in
$\mathbb F_p^2$.
\item Deduce that a binary quadratic form with non-zero discriminant takes every non-zero value of $\mathbb F_p$, and that
$x^2+y^2=1$ has $p-\chi(-1)$ solutions.
\end{enumerate}
\end{exercise}

\begin{exercise}[$\star$]\label{pb:hashing}
\emph{Universal hashing (Carter and Wegman, 1979).} Let $p$ be prime, $a,b$ independent and uniform in $\Z/p\Z$, and
$h(x)=\bigl((ax+b)\bmod p\bigr)\bmod m$ for keys $x\in\{0,\dots,p-1\}$, where $2\le m\le p$.
\begin{enumerate}[label=(\alph*)]
\item Show that for $x\ne y$ the pair $(ax+b,ay+b)$ is uniform on $(\Z/p\Z)^2$.
\item Deduce that $P(h(x)=h(y))\le\frac1m+\frac1p$ for $x\ne y$.
\item Store $n$ keys in $m=n^2$ buckets with such an $h$. Show that with probability at least $\frac12-\frac{n^2}{2p}$ no two keys
collide.
\end{enumerate}
\end{exercise}

\begin{exercise}[$\star\star$]\label{pb:triangleseigen}
\emph{Triangles from eigenvalues.} Let $G$ be a simple graph with $m$ edges, $t$ triangles and adjacency eigenvalues
$\lambda_1\ge\dots\ge\lambda_n$.
\begin{enumerate}[label=(\alph*)]
\item Show that $\sum_i\lambda_i^2=2m$ and $\sum_i\lambda_i^3=6t$.
\item Deduce that $t\le\frac{\sqrt2}3m^{3/2}$, and that the ratio $t/m^{3/2}$ tends to $\frac{\sqrt2}3$ for complete graphs.
\end{enumerate}
\end{exercise}

\begin{exercise}[$\star\star$]\label{pb:walks}
\emph{Many walks, many squares.} Let $G$ be a graph with $n$ vertices, $m$ edges and average degree $\bar d=\frac{2m}n$.
\begin{enumerate}[label=(\alph*)]
\item Show that for every even $k\ge2$ the number of walks of length $k$ in $G$ is at least $n\bar d^{\,k}$.
\item Show that $\operatorname{tr}A^4\ge\bar d^{\,4}$: the number of closed walks of length $4$ is at least what a random graph with the same
density has on average, $n^4(\bar d/n)^4$, up to lower-order terms.
\item For the Paley graph $P_p$ (Problem~\ref{pb:paley}), compute $\operatorname{tr}A^4$ and show that $\operatorname{tr}A^4=\bar d^{\,4}(1+O(\frac1p))$.
\end{enumerate}
\end{exercise}

\begin{exercise}[$\star\star$]\label{pb:littlewoodofford}
\emph{The Littlewood--Offord problem (Erdős, 1945).} Let $a_1,\dots,a_n$ be real numbers with $|a_i|\ge1$, and let $\varepsilon_1,\dots,\varepsilon_n$ be
independent uniform signs $\pm1$.
\begin{enumerate}[label=(\alph*)]
\item Show that for every open interval $J$ of length $2$, $P\bigl(\sum_i\varepsilon_ia_i\in J\bigr)\le\binom n{\lfloor n/2\rfloor}2^{-n}$.
\item Show that the bound is attained, and that it is asymptotic to $\sqrt{\frac2{\pi n}}$.
\end{enumerate}
\end{exercise}

\begin{exercise}[$\star\star$]\label{pb:cubeiso}
\emph{Edges inside a set of vertices of the cube.} Let $S$ be a set of vertices of the cube $\{0,1\}^n$, and $e(S)$ the number of edges of
the cube with both ends in $S$. Show that $e(S)\le\frac12|S|\log_2|S|$, with equality when $S$ is a subcube.
\end{exercise}

\begin{exercise}[$\star\star$]\label{pb:erdosselfridge}
\emph{The Erdős--Selfridge criterion (1973).} Maker and Breaker alternately claim unclaimed points of a finite set $V$, Maker first.
Maker wins if he claims all the points of some set of a given family $\mathcal H$ of subsets of $V$; otherwise Breaker wins.
\begin{enumerate}[label=(\alph*)]
\item Show that if $\sum_{H\in\mathcal H}2^{-|H|}<\frac12$, then Breaker has a winning strategy.
\item Maker and Breaker claim edges of the complete graph $K_N$, and Maker wants a complete subgraph $K_q$. Show that Breaker wins if
$\binom Nq2^{-\binom q2}<\frac12$, for instance if $N\le2^{q/2}$ and $q\ge3$.
\end{enumerate}
\end{exercise}

\begin{exercise}[Journey]\label{pb:shearer}
\emph{Shearer's lemma and the Loomis--Whitney inequality.} All random variables take finitely many values. For random variables
$X,Y$ define $H(Y\mid X)=\sum_xP(X=x)H(Y\mid X=x)$, where $H(Y\mid X=x)$ is the entropy of the conditional law.
\begin{enumerate}[label=(\alph*)]
\item Show the chain rule $H(X,Y)=H(X)+H(Y\mid X)$. Using subadditivity, $H(X,Y)\le H(X)+H(Y)$ (Volume~3, Exercise~4.9), deduce that
conditioning on more never raises entropy: $H(Y\mid X,Z)\le H(Y\mid X)$.
\item (Shearer, 1986.) Let $X=(X_1,\dots,X_n)$, and let $F_1,\dots,F_r\subseteq[n]$ be such that every $i\in[n]$ belongs to at least $k$ of them.
Writing $X_F=(X_i)_{i\in F}$, show that $kH(X)\le\sum_jH(X_{F_j})$.
\item (Loomis and Whitney, 1949, discrete form.) Let $S\subset\Z^3$ be finite, and $S_{xy},S_{yz},S_{xz}$ its projections on the coordinate
planes. Show that $|S|^2\le|S_{xy}|\,|S_{yz}|\,|S_{xz}|$.
\item Give a second proof of Problem~\ref{pb:triangleseigen}(b): a graph with $m$ edges has at most $\frac{(2m)^{3/2}}6$ triangles.
\end{enumerate}
\end{exercise}

\begin{exercise}[$\star\star$]\label{pb:sincpower}
\emph{Powers of $\frac{\sin t}t$.}
\begin{enumerate}[label=(\alph*)]
\item For $n\ge2$, show that
\[
\int_{-\infty}^\infty\Bigl(\frac{\sin t}t\Bigr)^n\dd t=\frac\pi{2^{n-1}(n-1)!}\sum_{0\le k<n/2}(-1)^k\binom nk(n-2k)^{n-1},
\]
and check the values $\pi,\frac{3\pi}4,\frac{2\pi}3$ for $n=2,3,4$.
\item Show that this integral is equivalent to $\sqrt{\frac{6\pi}n}$ as $n\to\infty$.
\end{enumerate}
\end{exercise}

\begin{exercise}[$\star\star$]\label{pb:cantor}
\emph{The Fourier transform of the Cantor law.} Let $B_1,B_2,\dots$ be independent fair bits and $X=\sum_{k\ge1}2B_k3^{-k}$, the
Cantor-distributed variable of Volume~3, Exercise~5.7. Let $\varphi(t)=\mathbb E[e^{itX}]$.
\begin{enumerate}[label=(\alph*)]
\item Show that $\varphi(t)=e^{it/2}\prod_{k\ge1}\cos\frac t{3^k}$.
\item Show that $|\varphi(2\pi3^m)|=|\varphi(2\pi)|$ for every $m\ge0$, and that $|\varphi(2\pi)|=0.3714\ldots>0$. Deduce, by a
Riemann--Lebesgue argument, that $X$ has no density.
\item By contrast, show that $\prod_{k\ge1}\cos\frac t{2^k}=\frac{\sin t}t$ (Viète's product), and deduce that $\sum_{k\ge1}\varepsilon_k2^{-k}$, with
independent fair signs $\varepsilon_k$, is uniform on $[-1,1]$.
\end{enumerate}
\end{exercise}

\begin{exercise}[$\star\star\star$]\label{pb:majoritystable}
\emph{Noise stability and Sheppard's formula.} For $-1<\rho<1$, a pair $(x,y)$ in $\{-1,1\}^n$ is \emph{$\rho$-correlated} if the pairs
$(x_i,y_i)$ are independent, with $x_i$, $y_i$ uniform and $\mathbb E[x_iy_i]=\rho$. The \emph{noise stability} of $f$ is
$\operatorname{Stab}_\rho(f)=\mathbb E[f(x)f(y)]$.
\begin{enumerate}[label=(\alph*)]
\item Show that $\operatorname{Stab}_\rho(f)=\sum_S\rho^{|S|}\hat f(S)^2$.
\item (Sheppard, 1899.) Let $(X,Y)$ be a centred Gaussian pair with unit variances and correlation $\rho$. Show that
$P(X>0,\,Y>0)=\frac14+\frac{\arcsin\rho}{2\pi}$, so that $\mathbb E[\operatorname{sgn}X\operatorname{sgn}Y]=\frac2\pi\arcsin\rho$.
\item Deduce that $\operatorname{Stab}_\rho(\operatorname{Maj}_n)\to\frac2\pi\arcsin\rho$ as $n\to\infty$ through odd values. (You may use the central limit theorem for
random vectors in $\R^2$, stated in Volume~3 after Theorem~6.15.)
\item Deduce that for every $k$, $W^k(\operatorname{Maj}_n)$ converges, to $0$ if $k$ is even and to $\frac2\pi\binom{2j}j\frac1{4^j(2j+1)}$ if $k=2j+1$; check
$W^1\to\frac2\pi$ against Problem~\ref{pb:influence}.
\end{enumerate}
\end{exercise}

\begin{exercise}[$\star\star\star$]\label{pb:period}
\emph{Cosets decide periodicity.} Let $\mu$ be a probability on a finite group $G$ whose support $S$ generates $G$, and let $X$ be the
random walk driven by $\mu$, started at $e$.
\begin{enumerate}[label=(\alph*)]
\item Show that the walk is irreducible.
\item Suppose that $S\subseteq gN$ for a normal subgroup $N\ne G$ and some $g\in G$. Show that $G/N$ is cyclic, generated by $gN$, and that the
walk has period at least $[G:N]$; in particular it does not converge to the uniform law.
\item Conversely, let $d$ be the period of the walk and $N$ the set of elements that the walk can reach at times divisible by
$d$. Show that $N$ is a normal subgroup of index $d$ and that $S\subseteq sN$ for any $s\in S$. Conclude: \emph{the walk converges to the
uniform law if and only if $S$ is contained in no coset of a proper normal subgroup.}
\item A random shuffle applies a uniformly chosen transposition to a deck of $n\ge2$ cards. Show that it never converges, and
compute its period.
\item Let $p\ge5$ be prime, and multiply a residue $x\in(\Z/p\Z)^*$ at each step by $2$ or by $3$, with probability $\frac12$ each. Assuming
that $2$ and $3$ generate $(\Z/p\Z)^*$, show that the walk converges to the uniform law if and only if $3\cdot2^{-1}$ is a primitive root
modulo $p$. Check that it fails for $p=5$ and $p=13$, and find the period in these two cases.
\end{enumerate}
\end{exercise}

\begin{exercise}[$\star\star$]\label{pb:rollingdie}
\emph{A die rolling on a chessboard.} A cube sits on a square of an infinite chessboard, exactly covering it. A move tips the cube
over one of the four edges of its bottom face, onto one of the four neighbouring squares.
\begin{enumerate}[label=(\alph*)]
\item Each move rotates the cube by a quarter turn about a horizontal axis. Using the action of the rotations on the four long
diagonals of the cube, show that every move is an odd permutation of the diagonals.
\item The cube makes a sequence of moves and comes back to its starting square with the same face on top. Show that it cannot
have turned by a quarter turn about the vertical axis. Can it have turned by a half turn?
\item Show that rolling the cube around the four squares of a $2\times2$ block (east, north, west, south) rotates it by a third
of a turn about a long diagonal. Deduce that every orientation allowed by (b), that is every even permutation of the
diagonals, can be obtained by a closed tour that comes back to the starting square.
\item Now the cube makes random moves, each of the four directions with probability $\frac14$. Show that the orientation of the cube
does not converge in law, but that the face on top does: after $t$ moves, each face is on top with probability tending to $\frac16$.
\item What is the expected number of moves before the cube is again in its initial orientation (wherever it is on the board)?
\end{enumerate}
\end{exercise}

\begin{exercise}[$\star\star\star$]\label{pb:gapdiam}
\emph{The diameter controls the speed.} Let $G$ be a finite group, $S$ a generating set with $S=S^{-1}$, and $d$ the diameter of the
Cayley graph: every element is a product of at most $d$ elements of $S$. Let $P(g,gs)=\frac1{|S|}$ for $s\in S$ be the symmetric walk.
For $f:G\to\R$ let
\[
\mathcal E(f)=\frac1{2|G||S|}\sum_{g\in G}\sum_{s\in S}\bigl(f(gs)-f(g)\bigr)^2,\qquad\operatorname{Var}(f)=\frac1{2|G|^2}\sum_{g,h\in G}\bigl(f(h)-f(g)\bigr)^2.
\]
\begin{enumerate}[label=(\alph*)]
\item Show that $\mathcal E(f)=\langle f,(I-P)f\rangle$ and $\operatorname{Var}(f)=\langle f,f\rangle-\langle f,\mathbf 1\rangle^2$ for the inner product $\langle f,h\rangle=\frac1{|G|}\sum_gf(g)h(g)$, and deduce
that $1-\lambda_2=\min\frac{\mathcal E(f)}{\operatorname{Var}(f)}$ over non-constant $f$, where $\lambda_2$ is the second largest eigenvalue of $P$.
\item Show that $\operatorname{Var}(f)\le|S|\,d^2\,\mathcal E(f)$ for every $f$. (Fix, for each $x\in G$, a word of length at most $d$ for $x$, and write $f(gx)-f(g)$ as a
telescoping sum along it.) Conclude that $1-\lambda_2\ge\frac1{|S|d^2}$.
\item For $\Z/n\Z$ with $S=\{\pm1\}$, compare with the true value $1-\cos\frac{2\pi}n$. For $\mathfrak S_n$ generated by the $n-1$ adjacent transpositions,
what lower bound do you get?
\end{enumerate}
\end{exercise}

\begin{exercise}[$\star\star\star$]\label{pb:treewalk}
\emph{A walk that escapes in a $4$-regular world.} The free group on $a,b$ is the set of reduced words in $a,a^{-1},b,b^{-1}$; its Cayley
graph is the infinite $4$-regular tree. A walk starts at the empty word $e$ and multiplies on the right by one of the four
letters, chosen uniformly. Let $D_t$ be the length of the reduced word at time $t$.
\begin{enumerate}[label=(\alph*)]
\item Show that $(D_t)$ is a Markov chain: from $0$ it goes to $1$, and from $k\ge1$ it goes to $k+1$ with probability $\frac34$ and to
$k-1$ with probability $\frac14$.
\item Show that the walk returns to $e$ with probability $\frac13$, and that the expected number of visits to $e$ (time $0$ included) is $\frac32$.
\item Let $h(k)=3^{-k/2}$. Show that $\mathbb E[h(D_{t+1})\mid D_t]\le\frac{\sqrt3}2h(D_t)$, and deduce that $P(D_{2n}=0)\le(\frac{\sqrt3}2)^{2n}$.
\item By counting the paths of $D$ that stay positive between times $1$ and $2n-1$, show that $P(D_{2n}=0)\ge\frac1{8n^2}(\frac34)^{n-1}$, and conclude that
$P(D_{2n}=0)^{1/2n}\to\frac{\sqrt3}2$.
\item Show that on the $q$-regular tree ($q\ge3$) the same limit is $\frac{2\sqrt{q-1}}q$. The simple walk on $\Z^2$ also has four neighbours at
every vertex: what is the limit there?
\end{enumerate}
\end{exercise}

\begin{exercise}[$\star\star\star$]\label{pb:commute}
\emph{The commute time identity.} Consider the simple random walk on $G$.
\begin{enumerate}[label=(\alph*)]
\item Let $a\in V$ and let $\tau\ge1$ be a stopping time with $X_\tau=a$ and $\mathbb E_a[\tau]<\infty$. Show that the expected number of visits to $x$ during
the times $0,1,\dots,\tau-1$ is $\pi(x)\,\mathbb E_a[\tau]$, for every $x\in V$.
\item Let $a\ne b$, and let $\tau$ be the first return to $a$ after the first visit to $b$. Using Volume~3, Exercise~7.8(b), show that the
walk visits $a$ on average $\deg(a)R(a,b)$ times before $\tau$, and deduce the commute time identity
\[
\mathbb E_a[T_b]+\mathbb E_b[T_a]=2|\mathcal E|\,R(a,b).
\]
\item Check the identity on a path of length $m$ between its two ends, and on a cycle of length $n$ between two neighbours.
\item Show that $\mathbb E_a[T_b]<2|\mathcal E|$ when $a$ and $b$ are neighbours, and deduce that the walk visits every vertex in expected time at most
$2|\mathcal E|(n-1)$, from any starting point. (Walk around a spanning tree.)
\end{enumerate}
\end{exercise}

\begin{exercise}[$\star\star\star$]\label{pb:foster}
\emph{Foster's theorem and symmetric graphs.}
\begin{enumerate}[label=(\alph*)]
\item Let $a\ne b$. Show that there is a vector $v\in\R^V$ with $Lv=e_a-e_b$, unique up to adding a constant, and that
$R(a,b)=v(a)-v(b)$. (Relate $v$ to the voltage $h$ of Volume~3, Exercise~7.8.)
\item Prove Foster's theorem: $\sum_{\{x,y\}\in\mathcal E}R(x,y)=n-1$. (Work on the hyperplane $W=\mathbf 1^\perp$, where $L$ is invertible, and take a trace.)
\item Suppose the automorphism group of $G$ acts transitively on the edges. Show that $R(x,y)=\frac{n-1}{|\mathcal E|}$ for every edge. If moreover, for every edge
$\{x,y\}$, some automorphism exchanges $x$ and $y$, show that $\mathbb E_x[T_y]=n-1$ for neighbours $x,y$.
\item Deduce the expected time for a walker on the vertices of a cube, of the Petersen graph, and of a dodecahedron to reach a
given neighbour of its starting vertex.
\end{enumerate}
\end{exercise}

\begin{exercise}[Journey]\label{pb:rayleigh}
\emph{Every subgraph of the plane grid is recurrent.} Let $a\ne b$ be vertices of $G$. A \emph{unit flow} from $a$ to $b$ is a function $\theta$ on
ordered pairs of neighbours with $\theta(x,y)=-\theta(y,x)$, $\sum_y\theta(x,y)=0$ for $x\ne a,b$, and $\sum_y\theta(a,y)=1$. Its \emph{energy} is
$\mathcal W(\theta)=\sum_{\{x,y\}\in\mathcal E}\theta(x,y)^2$.
\begin{enumerate}[label=(\alph*)]
\item With $v$ as in Problem~\ref{pb:foster}(a), show that $i(x,y)=v(x)-v(y)$ is a unit flow of energy $R(a,b)$, and that every unit flow has energy at least
$R(a,b)$ (Thomson's principle).
\item Deduce Rayleigh's monotonicity: removing edges from $G$ (keeping $a$ and $b$ connected) cannot decrease $R(a,b)$.
\item Let $\Pi_1,\dots,\Pi_m$ be disjoint sets of edges, each of which meets every path from $a$ to $b$. Show the Nash-Williams inequality
$R(a,b)\ge\sum_k\frac1{|\Pi_k|}$.
\item In $\Z^2$ let $B_n=\{-n,\dots,n\}^2$, and let $G_n$ be the graph obtained from $B_n$ (with the grid edges inside it) by merging its boundary into a
single vertex $z_n$. Show that $R(0,z_n)\ge\frac18\ln(2n+1)$.
\item Let $H$ be any set of edges of $\Z^2$ with at least one edge at $0$, and consider the walk that moves along the edges of $H$, to a
uniform $H$-neighbour. Show that it returns to $0$ with probability $1$.
\end{enumerate}
\end{exercise}

\begin{exercise}[Journey]\label{pb:aldousbroder}
\emph{A random walk draws a uniform spanning tree.} Run the simple random walk on $G$ from a vertex $r$ until it has visited every vertex,
and for each $v\ne r$ keep the edge along which the walk entered $v$ for the first time.
\begin{enumerate}[label=(\alph*)]
\item Show that the kept edges form a spanning tree, the \emph{first-entrance tree}.
\item An \emph{arborescence} rooted at $x$ is a spanning tree with every edge oriented towards $x$; each $u\ne x$ has one outgoing arc. The \emph{tree chain}
moves from $(T,x)$ as follows: pick a uniform neighbour $y$ of $x$, add the arc $x\to y$ and delete the arc leaving $y$. Show that the result is an
arborescence rooted at $y$. Show that if the tree chain follows the steps of a walk path that visits every vertex, its final state is the
\emph{last-exit tree} of the path, in which each vertex other than the endpoint points to the vertex visited just after its last visit.
\item Let $w(T,x)=\prod_{u\ne x}\frac1{\deg u}$. Show that $\mu(T,x)=\frac{w(T,x)}Z$ is a stationary law of the tree chain, that $\mu(T,x)=\frac{\deg x}{2|\mathcal E|\,\tau(G)}$ where $\tau(G)$
is the number of spanning trees, and that the second coordinate of a $\mu$-distributed state has law $\pi$.
\item Start the tree chain from $\mu$, run it for $N$ steps, and read the walk backwards (Volume~3, Proposition~8.14). Letting $N\to\infty$, prove the
theorem of Aldous and Broder: \emph{the first-entrance tree is uniformly distributed among the $\tau(G)$ spanning trees, whatever the starting
vertex $r$.}
\end{enumerate}
\end{exercise}

\begin{exercise}[$\star\star$]\label{pb:terras}
\emph{Collatz parities are fair coins.}
\begin{enumerate}[label=(\alph*)]
\item Show that $v_k(m)$ depends only on $m$ modulo $2^k$, and that $m\bmod2^k\mapsto v_k(m)$ is a bijection from $\Z/2^k\Z$ onto $\{0,1\}^k$.
\item Deduce that if $m$ is uniform on $2^k$ consecutive integers, the parities of $m,T(m),\dots,T^{k-1}(m)$ are independent fair coin
flips.
\item Show that $T^k(m)=\lambda m+\rho$, where $\lambda=\frac{3^j}{2^k}$, $j$ is the number of ones in $v_k(m)$, and $\rho\ge0$ depends only on $v_k(m)$.
\item Find the smallest $m\ge1$ whose first five parities are all odd. The number $1$ has parity vector $(1,0,1,0,1)$ of length $5$; find the
next integer with the same vector.
\end{enumerate}
\end{exercise}

\begin{exercise}[$\star\star\star$]\label{pb:everett}
\emph{Almost every integer goes down.} Let $\sigma(m)=\min\{k\ge1:T^k(m)<m\}$, with $\sigma(m)=\infty$ if there is no such $k$.
\begin{enumerate}[label=(\alph*)]
\item Let $k\ge1$ and let $u\in\{0,1\}^k$ have $j$ ones with $3^j<2^k$. Show that $\sigma(m)\le k$ for all large enough $m$ with $v_k(m)=u$.
\item Deduce that the set of $m$ with $\sigma(m)\le k$ has lower density at least $P\bigl(B_k<k\frac{\ln2}{\ln3}\bigr)$, where $B_k$ is binomial with parameters $k$
and $\frac12$.
\item Conclude (Terras, 1976; Everett, 1977): the set of $m$ with $\sigma(m)<\infty$ has natural density $1$.
\item In the model where the parities are independent fair coins for ever, show that $\ln T^k(m)$ drifts by $\frac12\ln\frac34$ per step on average, and
predict how many steps a typical $m$ needs to reach $1$. (A numerical calculation compares the prediction with the truth.)
\end{enumerate}
\end{exercise}

\begin{exercise}[Journey]\label{pb:coprimewalk}
\emph{A drunkard among the visible points.} A point $(x,y)\in\Z^2$ is \emph{visible} from the origin if $\gcd(x,y)=1$; a uniformly chosen point of a
large square is visible with probability close to $\frac6{\pi^2}$. Let $(X_t,Y_t)$ be the simple random walk on $\Z^2$ started at $0$.
\begin{enumerate}[label=(\alph*)]
\item Show that $U_t=X_t+Y_t$ and $V_t=X_t-Y_t$ are two independent simple random walks on $\Z$.
\item Let $p$ be an odd prime. Show that $P(p\mid X_t\text{ and }p\mid Y_t)=P(p\mid U_t)^2$, and that $P(p\mid U_t)\to\frac1p$. Show that $X_t$ and $Y_t$ are never both even when $t$ is
odd, and are both even with probability tending to $\frac12$ when $t$ is even.
\item Show that $P(p\mid U_t)\le\frac1p+\frac2{\sqrt t}$ for every odd prime $p$ and every $t\ge1$. (Write the probability as an average of $\cos^t\frac{2\pi k}p$.)
\item Prove that $P\bigl(\gcd(X_t,Y_t)=1\bigr)$ tends to $\frac8{\pi^2}$ along odd $t$, and to $\frac4{\pi^2}$ along even $t$. (For finitely many primes, use the walk of $(U_t,V_t)$
modulo $4\prod p$ and Problem~\ref{pb:period}; for the others, use (c).)
\end{enumerate}
\end{exercise}

\begin{exercise}[$\star\star$]\label{pb:range}
\emph{How much ground does a drunkard cover?} Let $S$ be the simple random walk on $\Z$ started at $0$, and $R_t$ the number of distinct
points among $S_0,\dots,S_t$. Write $u_{2k}=\binom{2k}k4^{-k}$.
\begin{enumerate}[label=(\alph*)]
\item Show that $\mathbb E[R_t]=\sum_{s=0}^tP(S_1\ne0,\dots,S_s\ne0)$. (A new point is found at time $s$ when $S_s$ differs from all earlier positions; read the path backwards
from time $s$.)
\item Using Volume~3, Theorem~7.3, show that $\mathbb E[R_{2m}]=(4m+1)\,u_{2m}$.
\item Deduce that $\mathbb E[R_t]\sim\sqrt{8t/\pi}$.
\item In $\Z^d$ with $d\ge3$, show that $\frac1t\mathbb E[R_t]$ tends to the probability that the walk never returns to $0$.
\end{enumerate}
\end{exercise}

\begin{exercise}[$\star\star$]\label{pb:cyclecover}
\emph{The last vertex of the ring.} A walker moves on the cycle $\Z/n\Z$ ($n\ge3$), one step clockwise or anticlockwise with probability
$\frac12$ each, starting at $0$.
\begin{enumerate}[label=(\alph*)]
\item Show that the set of visited vertices is always an arc containing $0$, and that once it has $k<n$ vertices, the expected
time to add a new one is $k$.
\item Deduce that the expected time to visit every vertex is $\frac{n(n-1)}2$.
\item Show that the last vertex to be visited is uniformly distributed on $\{1,\dots,n-1\}$.
\item Show that the same is true for the walk on the complete graph $K_n$, but not for the walk on the star with three branches
started at the end of a branch.
\end{enumerate}
\end{exercise}

\begin{exercise}[$\star\star$]\label{pb:threegamblers}
\emph{Three gamblers.} Three players start with $a$, $b$, $c$ coins ($a,b,c\ge1$), and $s=a+b+c$. In each round two of the players who still
have coins are chosen uniformly, and a fair coin decides which of the two gives one coin to the other. A player with no coins
leaves the game. Let $X_t$, $Y_t$, $Z_t$ be the fortunes after $t$ rounds.
\begin{enumerate}[label=(\alph*)]
\item Show that $X_tY_tZ_t+\frac s3t$ is a martingale as long as all three players are in the game, and deduce that the expected number of
rounds before the first player leaves is $\frac{3abc}s$.
\item Show that $X_tY_t+Y_tZ_t+Z_tX_t+t$ is a martingale until one player holds all the coins, and that the expected duration of the whole
game is $ab+bc+ca$.
\item Show that the first player ends up with all the coins with probability $\frac as$.
\item With $4$ coins each, how long is the game on average, and how long before the first player leaves?
\end{enumerate}
\end{exercise}

\begin{exercise}[$\star\star\star$]\label{pb:moran}
\emph{A small population forgets.} A population has $N$ individuals of two types $A$ and $B$. At each step one individual, chosen
uniformly, reproduces, and one individual, chosen uniformly and independently (possibly the same one), dies and is replaced by the
offspring. Let $X_t$ be the number of individuals of type $A$, with $X_0=i$. (This is the Moran model of population genetics.)
\begin{enumerate}[label=(\alph*)]
\item Show that $X$ is a martingale, and that type $A$ takes over the whole population with probability $\frac iN$.
\item Show that $\mathbb E[X_{t+1}(N-X_{t+1})\mid X_t]=(1-\frac2{N^2})X_t(N-X_t)$: the diversity decays geometrically.
\item Show that the expected time until one type has disappeared is
\[
t_i=Ni\sum_{x=i}^{N-1}\frac1x+N(N-i)\sum_{x=N-i+1}^{N-1}\frac1x.
\]
(The walk moves up or down with equal probability when it moves; count the expected number of moves made from each state.)
\item Deduce that a single mutant ($i=1$) is resolved in time $NH_{N-1}\sim N\ln N$, while for $i=\lfloor pN\rfloor$ with $0<p<1$,
\[
t_i\sim N^2\Bigl(p\ln\frac1p+(1-p)\ln\frac1{1-p}\Bigr).
\]
\end{enumerate}
\end{exercise}

\begin{exercise}[$\star\star\star$]\label{pb:bernoullilaplace}
\emph{Daniel Bernoulli's two urns.} Two urns contain $n$ balls each; altogether $n$ balls are red and $n$ are white. At each step a ball is
drawn uniformly from each urn, and the two balls are exchanged. Let $X_t$ be the number of red balls in the first urn.
\begin{enumerate}[label=(\alph*)]
\item Find the transition probabilities of $X$, and show that the chain is reversible with respect to the hypergeometric law
$\pi(x)=\binom nx^2\binom{2n}n^{-1}$.
\item Let $(Pf)(x)=\mathbb E[f(X_1)\mid X_0=x]$. Show that if $f$ is a polynomial of degree $k$, so is $Pf$, and that the coefficient of $x^k$ is multiplied by
\[
\lambda_k=1-\frac{k(2n-k+1)}{n^2}.
\]
Deduce that the eigenvalues of the transition matrix are $\lambda_0,\dots,\lambda_n$.
\item Start with all red balls in the first urn. Show that $\mathbb E[X_t]-\frac n2=\frac n2(1-\frac2n)^t$, and deduce that for $n\ge4$ and $t\le\frac n4\ln n$ this is at least
$\frac{\sqrt n}3$, while under $\pi$ the standard deviation of $X$ is less than $\frac{\sqrt n}2$.
\end{enumerate}
\end{exercise}

\begin{exercise}[$\star\star$]\label{pb:binompowers}
\emph{Powers of binomial coefficients.} Let $s\ge2$ be an integer. Show that
\[
\sum_{k=0}^n\binom nk^s\sim\frac{2^{sn}}{\sqrt s}\Bigl(\frac2{\pi n}\Bigr)^{(s-1)/2}.
\]
(Use the local de Moivre--Laplace theorem, Volume~3, Theorem~6.12, near the middle, and a cruder bound away from it.) Check the case $s=2$
against $\binom{2n}n$, and write the cases $s=3$ and $s=4$ explicitly.
\end{exercise}

\begin{exercise}[$\star\star$]\label{pb:saw}
\emph{Self-avoiding walks.} Let $c_n$ be the number of $n$-step walks on $\Z^2$ from the origin, with unit steps, that never visit a point twice. So $c_1=4$, $c_2=12$, $c_3=36$, $c_4=100$.
\begin{enumerate}[label=(\alph*)]
\item Show that $c_{m+n}\le c_mc_n$.
\item (Fekete's lemma.) Show that if a real sequence satisfies $a_{m+n}\le a_m+a_n$, then $\frac{a_n}n\to\inf_{k\ge1}\frac{a_k}k$ (which may be $-\infty$).
\item Deduce that $\mu=\lim c_n^{1/n}$ exists and equals $\inf_nc_n^{1/n}$.
\item Show that $1+\sqrt2\le\mu\le3$. (For the lower bound, count the walks that use only the steps east, north and south and never follow north by south or south by north.)
\item A numerical calculation counts $c_{14}=2374444$. What upper bound for $\mu$ does that give?
\end{enumerate}
\end{exercise}

\begin{exercise}[$\star\star$]\label{pb:inspection}
\emph{The inspection paradox for prime gaps.} Let $g_n=p_{n+1}-p_n$.
\begin{enumerate}[label=(\alph*)]
\item Show that $\sum_{p_{n+1}\le x}g_n^2\ge(1+o(1))\,x\log x$.
\item Let $p_N$ be the largest prime up to $x$, choose an integer $m$ uniformly in $]2,p_N]$, and let $G$ be the length of the gap $]p_n,p_{n+1}]$ containing $m$. Show that
$\mathbb E[G]=\frac1{p_N-2}\sum_{n<N}g_n^2$, and deduce that $\mathbb E[G]\ge(1+o(1))\log x$, although the average gap is $\sim\log x$. Compare with the inspection paradox
of Volume~3, Proposition~9.10.
\item In Cram\'er's model (Volume~5, Chapter~10) the gaps near $x$ are roughly geometric with mean $\log x$. Show that for a geometric law $G_0$ with mean $\frac1q$,
$\frac{\mathbb E[G_0^2]}{\mathbb E[G_0]}=\frac{2-q}q$, so that the model predicts $\sum_{p_{n+1}\le x}g_n^2\sim2x\log x$.
\item Figure~\ref{fig:gaps} shows $\frac1{x\log x}\sum_{p_{n+1}\le x}g_n^2$ for $10^3\le x\le10^8$. What do you conclude? If one draws $m$ at random and takes the next prime,
which primes are favoured?
\end{enumerate}
\end{exercise}

\begin{figure}[ht]
\centering
\fig{ch13-gaps}
\caption{The second moment of the prime gaps, $\frac1{x\log x}\sum_{p_{n+1}\le x}(p_{n+1}-p_n)^2$, for $x$ from $10^3$ to $10^8$ (blue): $1.16$ at $10^3$, $1.52$ at $10^6$,
$1.63$ at $10^8$. Problem~\ref{pb:inspection} proves that the limit inferior is at least $1$; Cram\'er's model predicts $2$.}
\label{fig:gaps}
\end{figure}

\begin{figure}[ht]
\centering
\fig{ch12-saw}
\caption{Self-avoiding walks on $\Z^2$: $c_n^{1/n}$ (blue) and $\frac{c_{n+1}}{c_n}$ (orange) for $n\le14$, with the bounds $1+\sqrt2$ and $3$ of Problem~\ref{pb:saw} and the
numerical value $\mu\approx2.638$ (dashed). Since $\mu=\inf_nc_n^{1/n}$, the blue curve always stays above $\mu$; it is still far from it at $n=14$.}
\label{fig:saw}
\end{figure}

\section{Hints}

\begin{hint}{ex:carowei}
$v$ is kept with probability $\frac1{d(v)+1}$: it must come first among itself and its $d(v)$ neighbours.
\end{hint}
\begin{hint}{ex:tournament}
Orient each game by a fair coin. A fixed set of $k$ players is beaten by a fixed outsider with probability $2^{-k}$.
\end{hint}
\begin{hint}{ex:sumfree}
The middle third is sum-free modulo $p$, and for each $a\ne0$, $ta\bmod p$ is uniform on $\{1,\dots,p-1\}$.
\end{hint}
\begin{hint}{ex:crossing}
Remove one edge per crossing. For the random subgraph, keep each vertex with probability $q=\frac{4n}m$ and take expectations of
$\operatorname{cr}\ge m-3n$ for the subgraph with the induced drawing.
\end{hint}
\begin{hint}{ex:propertyb}
A given set is monochromatic with probability $2^{1-k}$.
\end{hint}
\begin{hint}{ex:connectivity}
(a) $\mathbb E[Y]=n(1-p)^{n-1}$ and $\mathbb E[Y(Y-1)]=n(n-1)(1-p)^{2n-3}$. (b) Bound the expectation by $\sum_k\binom nkk^{k-2}p^{k-1}(1-p)^{k(n-k)}$ and use $1-p\le e^{-p}$. For the last
question, connectivity is preserved by adding edges.
\end{hint}
\begin{hint}{ex:giant}
(b) Chernoff. (c) At that moment $|U|\le1+\varepsilon n$ and all pairs between $S$ and $T$ have been queried. (d) If $|U|<\frac{\varepsilon^2n}5$, then $|S|$ is large, and
$|S|\cdot|T|$ exceeds $N$.
\end{hint}

\begin{hint}{pb:vanish}
(a) $\Phi_p$ is irreducible. (b) Write everything in the basis $1,\zeta$ of $\Q(\zeta_6)$, using $\zeta^3=-1$. (c) Group the
exponents by their residue modulo $p^{a-1}$; over $\Q(\omega)$, $\zeta_{p^a}$ has degree $p^{a-1}$.
\end{hint}

\begin{hint}{pb:samecycle}
(a) Count the permutations in which the cycle of $1$ has length $\ell$. (b) Condition on that length.
\end{hint}

\begin{hint}{pb:denes}
Multiplying by a transposition $(a\,b)$ merges two cycles if $a$ and $b$ are in different cycles, and splits one otherwise.
\end{hint}

\begin{hint}{pb:conics}
(a) Complete the square: $4a(ax^2+bx+c)=(2ax+b)^2-D$. Then count the solutions of $y^2-z^2=D$. (b) For fixed $y$, the number of
$x$ is $1+\chi(\text{discriminant})$.
\end{hint}

\begin{hint}{pb:hashing}
(c) Expected number of colliding pairs, and Markov's inequality.
\end{hint}

\begin{hint}{pb:triangleseigen}
(b) $\sum\lambda_i^3\le\lambda_1\sum\lambda_i^2$ and $\lambda_1^2\le\sum\lambda_i^2$.
\end{hint}

\begin{hint}{pb:walks}
(a) Expand $\mathbf1$ in an orthonormal basis of eigenvectors and use Jensen. (b) $\lambda_1\ge\bar d$ (Rayleigh). (c) Use the counts of
common neighbours in Problem~\ref{pb:paley}(a).
\end{hint}

\begin{hint}{pb:littlewoodofford}
Make all $a_i\ge1$, then show that the sets $\{i:\varepsilon_i=1\}$ giving sums in $J$ form an antichain.
\end{hint}

\begin{hint}{pb:cubeiso}
Induction on $n$; split $S$ by the last coordinate. You will need $t\log_2t+2t\le(1+t)\log_2(1+t)$ on $[0,1]$.
\end{hint}

\begin{hint}{pb:erdosselfridge}
Give each set not yet touched by Breaker the weight $2^{-r}$, where $r$ is the number of its points Maker has not yet claimed, and let Breaker always reduce the
total weight as much as possible.
\end{hint}

\begin{hint}{pb:shearer}
(a) Write $\log\frac1{P(X=x,Y=y)}$ as a sum of two logarithms. (b) Chain rule, then drop part of the conditioning. (d) A random ordered triangle.
\end{hint}

\begin{hint}{pb:sincpower}
(a) $\frac{\sin t}t$ is the characteristic function of the uniform law on $[-1,1]$; invert at $0$. (b) Laplace's method on $[-\pi,\pi]$,
where $\log\frac{\sin t}t=-\frac{t^2}6+O(t^4)$.
\end{hint}

\begin{hint}{pb:cantor}
(a) Independence turns the characteristic function of a sum into a product. (b) Multiplying $t$ by $3$ shifts the product by one
factor, and the extra factor is $\cos(2\pi\cdot\text{integer})$. For Riemann--Lebesgue, approximate a density by step functions.
(c) $\sin t=2^m\sin\frac t{2^m}\prod_{k=1}^m\cos\frac t{2^k}$.
\end{hint}

\begin{hint}{pb:majoritystable}
(a) $\mathbb E[x^Sy^T]$. (b) Write $Y=\rho X+\sqrt{1-\rho^2}Z$ and use the rotation invariance of $(X,Z)$: the event is a wedge. (c) $(\frac{\sum x_i}{\sqrt n},\frac{\sum y_i}{\sqrt n})$.
(d) Power series with non-negative coefficients of sum $1$ that converge pointwise on $[0,1[$; the Taylor series of $\arcsin$.
\end{hint}

\begin{hint}{pb:period}
(a) In a finite group the semigroup generated by $S$ is a group. (b) Work in $G/N$, where the walk moves deterministically. (c) The cyclic classes of
Volume~3, Lemma~8.4 are $N,sN,s^2N,\dots$; to see that $N$ is normal, write $s^{-1}$ as a product of elements of $S$. (d) Signs. (e) $S=\{2,3\}\subseteq2\langle3\cdot2^{-1}\rangle$.
\end{hint}

\begin{hint}{pb:rollingdie}
(a) Volume~2, Proposition~1.5 and Exercise~1.6: a quarter turn about a face axis is a $4$-cycle of the diagonals. (b) Colour the board like a
chessboard. (c) A rotation of order $3$ fixes a diagonal; two such rotations about different diagonals generate $A_4$. (d) Problem~\ref{pb:period}, and
the action of $A_4$ on the six faces. (e) The uniform law is stationary.
\end{hint}

\begin{hint}{pb:gapdiam}
(a) $P$ is a symmetric matrix. (b) For $x=s_1\cdots s_\ell$, $(f(gx)-f(g))^2\le\ell\sum_i(f(gs_1\cdots s_i)-f(gs_1\cdots s_{i-1}))^2$; then sum over $g$, and substitute $g'=gs_1\cdots s_{i-1}$.
(c) The diameters are $\lfloor\frac n2\rfloor$ and $\binom n2$.
\end{hint}

\begin{hint}{pb:treewalk}
(b) Volume~3, Corollary~7.10. (c) Compute $\frac14h(k-1)+\frac34h(k+1)$. (d) Catalan numbers, and $\binom{2m}m\ge\frac{4^m}{2m+1}$. (e) The same computations with $\frac1q$ and $\frac{q-1}q$.
\end{hint}

\begin{hint}{pb:commute}
(a) Count visits along the cycles of length $\tau$ and use the strong Markov property: $\nu(x)=\mathbb E_a[\text{visits to }x\text{ before }\tau]$ satisfies $\nu P=\nu$. (b) Each
excursion from $a$ reaches $b$ before coming back with probability $\frac1{\deg(a)R}$. (d) $R(a,b)\le1$ along an edge; a closed walk round a
spanning tree uses each tree edge once in each direction.
\end{hint}

\begin{hint}{pb:foster}
(a) $L$ is symmetric with kernel $\R\mathbf1$, so its image is $W$. Take $v=-\frac hI$. (b) $R(x,y)=\langle e_x-e_y,L^{-1}(e_x-e_y)\rangle$ on $W$, and $L=\sum_{\{x,y\}}(e_x-e_y)(e_x-e_y)^{\mathsf T}$. (c) An
automorphism preserves resistances; the exchange gives $\mathbb E_x[T_y]=\mathbb E_y[T_x]$.
\end{hint}

\begin{hint}{pb:rayleigh}
(a) For a unit flow $\theta$, write $\theta=i+\eta$ and show $\sum_{\{x,y\}}i(x,y)\eta(x,y)=0$ by summation by parts. (b) A flow on the smaller graph is a flow on
the larger one. (c) The flow through each $\Pi_k$ is at least $1$; Cauchy--Schwarz. (d) Cut between the boundaries of $B_k$ and $B_{k+1}$. (e) Volume~3,
Exercise~7.8(b) on $H\cap B_n$ with its boundary merged.
\end{hint}

\begin{hint}{pb:aldousbroder}
(b) Each step resets the arc of the vertex that is left. (c) For a given $(T',y)$, the predecessors are obtained by choosing a child $x$ of $y$ and a
neighbour $z$ of $y$ in the subtree of $x$. (d) The last-exit tree of the forward path is the first-entrance tree of the reversed path.
\end{hint}

\begin{hint}{pb:terras}
(a) Induction on $k$: $T$ maps the two classes $r$ and $r+2^k$ modulo $2^{k+1}$ to classes modulo $2^k$ that differ by $3^{\epsilon}2^{k-1}$. (c) Induction.
\end{hint}

\begin{hint}{pb:everett}
(a) $T^k(m)=\lambda m+\rho$ with $\lambda<1$. (b) Problem~\ref{pb:terras}(a): the vectors with few ones correspond to that many residue classes. (c) $\frac{\ln2}{\ln3}>\frac12$ and the
weak law of large numbers. (d) Each step multiplies by $\frac12$ or roughly $\frac32$.
\end{hint}

\begin{hint}{pb:coprimewalk}
(a) Each step changes $U$ and $V$ by $\pm1$, with the four sign pairs equally likely. (b) $X=\frac{U+V}2$ and $Y=\frac{U-V}2$. (c) $|\cos y|\le e^{-2y^2/\pi^2}$ for $|y|\le\frac\pi2$, and
a Gaussian integral. (d) Truncate at primes $\le P$ and let $P\to\infty$ after $t$.
\end{hint}

\begin{hint}{pb:range}
(a) Reverse the steps: $S_s\ne S_r$ for all $r<s$ iff the reversed walk does not return to $0$ in $s$ steps. (b) Returns only happen at even times, and
$\sum_{k=0}^mu_{2k}=(2m+1)u_{2m}$. (c) $u_{2m}\sim\frac1{\sqrt{\pi m}}$. (d) Cesàro.
\end{hint}

\begin{hint}{pb:cyclecover}
(a) Gambler's ruin on an interval of length $k+1$. (c) For $v$ to be last, the walk must reach one neighbour of $v$ and then go all the way round
to the other neighbour without touching $v$. (d) On the star, where can the walk be last?
\end{hint}

\begin{hint}{pb:threegamblers}
(a) Compute the conditional expectation of the change for each pair, then average. Optional stopping needs $\mathbb E[\tau]<\infty$: Volume~3, Lemma~7.8 on a finite
state space. (c) $X_t$ is a bounded martingale.
\end{hint}

\begin{hint}{pb:moran}
(b) $\mathbb E[(X_{t+1}-X_t)^2\mid X_t]=\frac{2X_t(N-X_t)}{N^2}$. (c) From $x$, a move happens with probability $q_x=\frac{2x(N-x)}{N^2}$, so each move from $x$ costs $\frac1{q_x}$ steps on
average; the simple walk on $\{0,\dots,N\}$ started at $i$ makes on average $\frac{2i(N-x)}N$ moves from $x\ge i$, and $\frac{2x(N-i)}N$ from $x\le i$. (d) Riemann sums.
\end{hint}

\begin{hint}{pb:bernoullilaplace}
(a) Detailed balance. (b) Expand $(x+\delta)^k$ where $\delta\in\{-1,0,1\}$ is the change: only $\mathbb E[\delta]$ and $\mathbb E[\delta^2]$ reach degree $k$. (c) $\ln(1-y)\ge-y-y^2$ for $0\le y\le\frac12$.
\end{hint}

\begin{hint}{pb:binompowers}
With $p_k=\binom nk2^{-n}$, the sum is $2^{sn}\sum_kp_k^s$. Near the middle, $p_k\approx\frac2{\sqrt n}\phi(x_k)$ with $x_k=\frac{2k-n}{\sqrt n}$, and the $x_k$ are spaced $\frac2{\sqrt n}$ apart.
Away from it, $\sum p_k^s\le(\max_kp_k)^{s-1}P(|x_K|>A)$. Finally $\int_\R\phi^s=(2\pi)^{(1-s)/2}s^{-1/2}$.
\end{hint}

\begin{hint}{pb:saw}
(a) Cut a walk of length $m+n$ after $m$ steps. (b) Fix $k$, write $n=qk+r$. (d) After the first step, a self-avoiding walk never goes straight back. For the lower bound, set up a
two-state recurrence (last step horizontal or vertical).
\end{hint}

\begin{hint}{pb:inspection}
(a) Cauchy--Schwarz, $\sum g_n\le x$ with $\pi(x)$ terms. (b) Each $m$ in a gap of length $g$ sees $G=g$, and there are $g$ such $m$. (c) $\mathbb E[G_0]=\frac1q$, $\mathbb E[G_0^2]=\frac{2-q}{q^2}$.
\end{hint}

\section{Solutions}

\begin{solution}{ex:carowei}
\idea A random order resolves every edge in favour of one endpoint, and a vertex wins against all its neighbours with
probability one over the size of its neighbourhood.
\writeup If two kept vertices were adjacent, each would come before the other. Vertex $v$ is kept when it is first among the $d(v)+1$
vertices of its closed neighbourhood, which happens with probability $\frac1{d(v)+1}$ by symmetry. By linearity the expected size of the kept set is
$\sum_v\frac1{d(v)+1}$, and some order achieves at least this. By convexity of $x\mapsto\frac1{x+1}$ (Jensen, Theorem~\ref{thm:jensenprob} for the uniform law on the vertices), this is at
least $\frac n{d+1}$, where $d$ is the average degree: about twice the bound $\frac n{2d}$ of Proposition~\ref{prop:alteration}(b). For a disjoint union of cliques $K_{d+1}$ it is exact.
\end{solution}

\begin{solution}{ex:tournament}
\idea Count the bad events: a set of $k$ players that nobody beats entirely. Each outsider fails independently.
\writeup Orient each game by an independent fair coin. For a set $K$ of $k$ players and an outsider $w$, $w$ beats all of $K$ with probability $2^{-k}$, and
these events are independent for the $n-k$ outsiders, which involve disjoint sets of games. So $K$ is unbeaten with probability $(1-2^{-k})^{n-k}$,
and the expected number of unbeaten sets is $\binom nk(1-2^{-k})^{n-k}<1$: some tournament has none. Numerically, $\binom{21}2(\frac34)^{19}\approx0.89$ and $\binom{91}3(\frac78)^{88}\approx0.96$. (The smallest
tournament for $k=2$ has $7$ players: let $i$ beat $j$ when $j-i$ is a non-zero square modulo $7$.)
\end{solution}

\begin{solution}{ex:sumfree}
\idea Multiplying by a random unit spreads each element uniformly over the non-zero residues, and the middle third is a
sum-free target of relative size just above one third.
\writeup Let $M=\{q+1,\dots,2q+1\}$, which has $q+1$ elements. If $x,y\in M$ then $2q+2\le x+y\le4q+2$, and modulo $p=3q+2$ the sum lies in $\{2q+2,\dots,3q+1\}\cup\{0,\dots,q\}$, outside $M$:
$M$ is sum-free modulo $p$. For $a\in A$, $a\not\equiv0\pmod p$ since $p>|a|$, so $t\mapsto ta\bmod p$ is a bijection of $\{1,\dots,p-1\}$ ($\Z/p\Z$ is a field), and $ta\bmod p$ is uniform: it lies in
$M$ with probability $\frac{q+1}{3q+1}>\frac13$. So the expected number of kept elements exceeds $\frac n3$, and some $t$ keeps more than $\frac n3$. The kept set $A_t$ is
sum-free: if $x+y=z$ in $A_t$, then $tx+ty\equiv tz\pmod p$ with all three residues in $M$, which is impossible. (Infinitely many primes $p\equiv2\pmod3$ exist, by
the argument of Euclid applied to $3N!-1$.)
\end{solution}

\begin{solution}{ex:crossing}
\idea A drawing with few crossings would give random subgraphs with fewer crossings than their many edges force.
\writeup Take a drawing with $\operatorname{cr}(G)$ crossings and delete one edge from each crossing: the remaining graph is drawn without crossings, so it is
planar and has at most $3n$ edges; hence $m-\operatorname{cr}(G)\le3n$ (also trivially true for $n\le2$). Now keep each vertex with probability $q$, independently, and consider the
induced subgraph $G'$ with the induced drawing, with $n'$ vertices, $m'$ edges and $c'$ crossings in that drawing. A crossing involves two edges with four
distinct endpoints (in an optimal drawing, edges with a common endpoint do not cross), so $\mathbb E[c']=q^4\operatorname{cr}(G)$, while $\mathbb E[m']=q^2m$ and $\mathbb E[n']=qn$. As $c'\ge\operatorname{cr}(G')\ge m'-3n'$ for
every outcome, $q^4\operatorname{cr}(G)\ge q^2m-3qn$. With $q=\frac{4n}m\le1$:
\[
\operatorname{cr}(G)\ge\frac m{q^2}-\frac{3n}{q^3}=\frac{m^3}{16n^2}-\frac{3m^3}{64n^2}=\frac{m^3}{64n^2}.
\]
This short proof, from the book of Aigner and Ziegler, has applications from incidence geometry (the Szemer\'edi--Trotter
theorem) to sum-product estimates.
\end{solution}

\begin{solution}{ex:propertyb}
\idea Exactly as for Ramsey numbers and for $k$-SAT: a random colouring, then the union bound or the local lemma.
\writeup Colour each element red or blue at random. A given set is monochromatic with probability $2\cdot2^{-k}=2^{1-k}$. With fewer than $2^{k-1}$ sets, the expected
number of monochromatic sets is less than $1$. In the second case, the event that a set is monochromatic is mutually independent of the
events for the sets disjoint from it (Proposition~\ref{prop:indepfun}), and Theorem~\ref{thm:lll} with $p=2^{1-k}$ applies.
\end{solution}

\begin{solution}{ex:connectivity}
\idea Below the threshold the graph still has isolated vertices; above it, the only way to be disconnected would be a
small piece cut off, and these are too unlikely.
\writeup (a) $\mathbb E[Y]=n(1-p)^{n-1}$. Since $\ln(1-p)=-p+O(p^2)$ and $np^2\to0$, $(1-p)^{n-1}=e^{-np+o(1)}=\frac{e^{-\omega}}ne^{o(1)}$, so $\mathbb E[Y]\sim e^{-\omega}$ (when $p\to0$, which holds if $\omega=o(\ln n)$).
Two given vertices are both isolated with probability $(1-p)^{2n-3}$ (they avoid $2n-3$ edges), so $\mathbb E[Y^2]=\mathbb E[Y]+n(n-1)(1-p)^{2n-3}$ and
\[
\begin{aligned}\frac{\operatorname{Var}(Y)}{\mathbb E[Y]^2}&\le\frac1{\mathbb E[Y]}+\frac{(1-p)^{-1}n(n-1)}{n^2}-1\eqbr\le\frac1{\mathbb E[Y]}+\frac p{1-p}.\end{aligned}
\]
If $\omega\to-\infty$ (with $\omega\ge-\ln n$, so that $p\ge0$), then $p\to0$ and, using $\ln(1-p)\ge-p-p^2$ for $p\le\frac12$,
$\mathbb E[Y]\ge n(1-p)^n\ge e^{\ln n-np-np^2}=e^{-\omega-np^2}\to\infty$, so $P(Y=0)\to0$ by Lemma~\ref{lem:secondmoment}: w.h.p.\ there is an isolated vertex, and the graph is disconnected.

(b) If the graph is disconnected, some component $S$ has $1\le|S|\le\frac n2$ vertices: it is connected, so contains one of the $k^{k-2}$ spanning trees on $S$, whose
$k-1$ edges are present, and the $k(n-k)$ pairs between $S$ and its complement are absent. The expected number of such pairs (set, tree) is
\[
E=\sum_{k=1}^{n/2}\binom nkk^{k-2}p^{k-1}(1-p)^{k(n-k)}.
\]
Using $\binom nk\le\frac{n^k}{k!}\le\bigl(\frac{en}k\bigr)^k$, $k^{k-2}\le k^k$, $(1-p)^{k(n-k)}\le e^{-pk(n-k)}$ the $k$-th term is at most $\frac1p\bigl(enp\,e^{-p(n-k)}\bigr)^k$, that is
\[
\frac1p\Bigl(e(\ln n+\omega)\,e^{-(\ln n+\omega)(1-\frac kn)}\Bigr)^k.
\]
For $k\le\frac n{\ln n}$, $(\ln n+\omega)\frac kn\le2$, so the base is at most $\frac{e^3(\ln n+\omega)e^{-\omega}}n$ and the sum over these $k$ is at most $\frac1p\cdot2\frac{e^3(\ln n+\omega)e^{-\omega}}n\le2e^3e^{-\omega}\to0$
(a geometric series with ratio $\le\frac12$ for large $n$, and $\frac1p\cdot\frac{\ln n+\omega}n=1$). For $\frac n{\ln n}<k\le\frac n2$, the base is at most $e(\ln n+\omega)e^{-(\ln n+\omega)/2}\le\frac{e\cdot2\ln n}{\sqrt n}$, which is less
than $n^{-1/3}$ for large $n$, so these terms add up to at most $\frac n{p}\,n^{-\frac{n}{3\ln n}}\to0$. So $E\to0$, and by Markov's inequality the graph is connected w.h.p.
Finally, if $\omega\to\infty$ faster, replace $\omega$ by $\min(\omega,\ln\ln n)$: the graph $G(n,p)$ can be obtained from $G(n,p')$ with $p'\le p$ by adding
independent edges (keep each absent edge with probability $\frac{p-p'}{1-p'}$), and adding edges preserves connectivity. Together with (a), $\frac{\ln n}n$ is a
sharp threshold: with $\omega=c$ constant, $P(\text{connected})\to e^{-e^{-c}}$, the probability that a Poisson$(e^{-c})$ number of isolated vertices is zero (Erd\H os and R\'enyi, 1959).
\end{solution}

\begin{solution}{ex:giant}
\idea Depth-first search keeps its current branch on a stack, which is a path. If the stack stayed short, many
vertices would be finished early, and all the pairs between them and the unvisited vertices would have been queried
with no success: too many queries for the time available.
\writeup (a) A vertex is pushed onto $U$ when it is adjacent to the current top, so $U$ is always a path, bottom to top, and lies in one component.
When $v$ moves to $S$, all pairs $\{v,w\}$ with $w$ in $T$ at that moment have been queried, negatively; $T$ only shrinks, so all pairs between $S$ and $T$ have
been queried. Each positive answer moves a new vertex from $T$ to $U$, and vertices never return to $T$, so $|S|+|U|$ is at least the number of positive
answers. Moreover $|U|$ exceeds by at most $1$ the number of positive answers since $U$ was last empty.

(b) The number $P_N$ of positive answers among the first $N$ is binomial$(N,p)$, with mean $Np=\frac{\varepsilon(1+\varepsilon)n}2+O(1)$. By Chernoff (Theorem~\ref{thm:chernoff}) with $\delta=n^{2/3}/(Np)$,
$P(P_N\le Np-n^{2/3})\le e^{-n^{4/3}/(2Np)}\to0$, since $Np$ is of order $n$.

(c) The search makes at most $\binom n2>N$ queries before it ends, so it is still running at time $N$. Suppose $|S|\ge\frac n3$ at time $N$; consider the first moment
when $|S|=\lceil\frac n3\rceil$. Then $|U|\le1+P_N$ by (a), and w.h.p.\ $P_N\le2Np\le1.1\varepsilon n$ (Chernoff again), so
$|T|\ge n-\frac n3-1-1.1\varepsilon n\ge\frac n2$ for large $n$ (as $\varepsilon\le\frac1{10}$). All $|S|\,|T|\ge\frac{n^2}6$ pairs between $S$ and $T$ have been queried, but $\frac{n^2}6>\frac{\varepsilon n^2}2\ge N$: a contradiction.

(d) Suppose $|U|<\frac{\varepsilon^2n}5$ at time $N$. By (a) and (b), $|S|\ge P_N-|U|\ge\frac{\varepsilon n}2+\frac{3\varepsilon^2n}{10}-n^{2/3}$, and by (c) $|S|<\frac n3$. Then
$|T|=n-|S|-|U|$, and since $s\mapsto s(n-s-|U|)$ increases for $s\le\frac{n-|U|}2$,
\[
\begin{gathered}|S|\,|T|\ge s_0\Bigl(n-s_0-\frac{\varepsilon^2n}5\Bigr),\qq s_0=\frac{\varepsilon n}2+\frac{3\varepsilon^2n}{10}-n^{2/3},\end{gathered}
\]
which equals $n^2\bigl(\frac\varepsilon2+\frac{\varepsilon^2}{20}-\frac{2\varepsilon^3}5-\frac{3\varepsilon^4}{20}\bigr)+O(n^{5/3})$. For $\varepsilon\le\frac1{10}$, $\frac{\varepsilon^2}{20}>\frac{2\varepsilon^3}5+\frac{3\varepsilon^4}{20}$, so for large $n$ this
exceeds $\frac{\varepsilon n^2}2\ge N$, although all these pairs have been queried: a contradiction. So w.h.p.\ $|U|\ge\frac{\varepsilon^2n}5$, and $U$ is a path in one component. For
$G(n,\frac cn)$ with $c>1$, take $\varepsilon=\min(c-1,\frac1{10})$: $G(n,\frac cn)$ contains a copy of $G(n,\frac{1+\varepsilon}n)$ (add independent edges as in Exercise~\ref{ex:connectivity}), and adding edges
does not shrink components.
\end{solution}

\begin{solution}{pb:vanish}
\idea Relations between roots of unity are polynomials divisible by the minimal polynomial. For a prime $p$, the minimal
polynomial $\Phi_p$ has all coefficients $1$, which forces the coefficients to be equal. For $n=6$, a basis of
$\Q(\zeta_6)$ turns the question into two linear equations in $0$/$1$ unknowns. For $n=p^a$, $\zeta_{p^a}$ generates
$\Q(\zeta_{p^a})$ over $\Q(\omega)$ with a binomial minimal polynomial, so a relation splits into relations among $p$-th roots.

\writeup (a) Let $A=\sum_{k<p}a_kX^k\in\Q[X]$, of degree $<p$. If $A(\zeta)=0$, the minimal polynomial $\Phi_p$ of $\zeta$
(irreducible, Volume~0, Chapter~3) divides $A$; by degrees $A=c\Phi_p$, so all $a_k=c$. Conversely
$c(1+\zeta+\dots+\zeta^{p-1})=0$. For sets ($a_k\in\{0,1\}$) the only vanishing sums are the empty set and all $p$ roots.

(b) Let $\zeta=e^{i\pi/3}$, so $\zeta^3=-1$, and let $c_k\in\{0,1\}$ indicate whether $\zeta^k$ is in the set. Put
$u_k=c_k-c_{k+3}$ for $k=0,1,2$: the sum is $u_0+u_1\zeta+u_2\zeta^2$. As $\zeta^2=\zeta-1$, this is
$(u_0-u_2)+(u_1+u_2)\zeta$, and since $1,\zeta$ are linearly independent over $\R$, the sum vanishes if and only if
$u_0=u_2=-u_1=t$. If $t=0$, the set is a union of antipodal pairs $\{\zeta^k,-\zeta^k\}$: $7$ non-empty sets. If $t=1$,
$c_0=c_2=c_4=1$ and $c_1=c_3=c_5=0$: the triangle $\{1,\zeta^2,\zeta^4\}$. If $t=-1$, the other triangle. That makes $9$.

(c) Let $\zeta=e^{2i\pi/p^a}$, $N=p^{a-1}$ and $\omega=\zeta^N$, a primitive $p$-th root. By the tower law and
$[\Q(\zeta):\Q]=\varphi(p^a)$ (Volume~2, Chapter~10), $[\Q(\zeta):\Q(\omega)]=\frac{\varphi(p^a)}{p-1}=N$, and $\zeta$ is a root
of $X^N-\omega$, so that is its minimal polynomial over $\Q(\omega)$ and $1,\zeta,\dots,\zeta^{N-1}$ is a basis of $\Q(\zeta)$
over $\Q(\omega)$ (Volume~2, Chapter~8). Write each exponent as $r+Nj$, $0\le r<N$, $0\le j<p$. A set with indicator $c$ has
sum $\sum_{r<N}\zeta^r\bigl(\sum_{j<p}c_{r+Nj}\omega^j\bigr)$. It vanishes if and only if each inner sum vanishes, and by (a)
this means that for each $r$, either all $\zeta^{r+Nj}=\zeta^r\omega^j$ are in the set or none. These are the sets
$\{\eta,\eta\omega,\dots,\eta\omega^{p-1}\}$ with $\eta=\zeta^r$.
\tools Volume~0, Chapter~3 (irreducibility of $\Phi_p$); Volume~2, Chapters~8 and~10 (tower law, cyclotomic fields).
\leads For $n$ with two prime factors, (b) generalises: every vanishing sum of $n$-th roots with non-negative coefficients
is a sum of rotated regular $p$-gons and $q$-gons (de Bruijn, 1953), but with three primes, as for $\Phi_{105}$, this
fails, and cancellations between polygons appear. Lam and Leung (2000) proved that a vanishing sum of $m$ $n$-th roots of
unity, repetitions allowed, exists if and only if $m$ is a combination with non-negative integer coefficients of the
primes dividing $n$.
\end{solution}

\begin{solution}{pb:samecycle}
\idea Building the cycle of $1$ step by step: it has length $\ell$ for $(n-1)(n-2)\cdots(n-\ell+1)\cdot(n-\ell)!=(n-1)!$
permutations, whatever $\ell$ is. Then, given the length, the other points of the cycle are a uniform random subset.

\writeup (a) A permutation in which the cycle of $1$ has length $\ell$ is given by the sequence $\sigma(1),\sigma^2(1),\dots,\sigma^{\ell-1}(1)$
of distinct elements different from $1$, in $(n-1)\cdots(n-\ell+1)$ ways, and a permutation of the remaining $n-\ell$ points, in
$(n-\ell)!$ ways. That is $(n-1)!$, so the length is uniform.

(b) Given that the cycle of $1$ has length $\ell$, by symmetry its other $\ell-1$ points are a uniform $(\ell-1)$-subset of
$\{2,\dots,n\}$, which contains $\{2,\dots,k\}$ with probability $\binom{n-k}{\ell-k}/\binom{n-1}{\ell-1}=\binom{\ell-1}{k-1}/\binom{n-1}{k-1}$.
Averaging over $\ell$ and using $\sum_{\ell=k}^n\binom{\ell-1}{k-1}=\binom nk$:
\[
P=\frac1n\sum_{\ell=1}^n\frac{\binom{\ell-1}{k-1}}{\binom{n-1}{k-1}}=\frac{\binom nk}{n\binom{n-1}{k-1}}=\frac1k.
\]
\tools Volume~3, Chapters~1 and~2 (conditioning); Volume~0, Chapter~1 (the hockey-stick identity).
\leads The answer does not depend on $n$, a sign of the \emph{Chinese restaurant process}: customers $1,2,\dots$ arrive, and each
joins an existing cycle or opens a new one, which builds a random permutation one point at a time (Volume~3, Theorem~2.12).
\end{solution}

\begin{solution}{pb:denes}
\idea Follow the cycles as the transpositions are multiplied one at a time. Starting from $n$ cycles, each step changes the number
of cycles by $\pm1$, so reaching one cycle in $n-1$ steps requires a merge every time, and merges are exactly the edges that join two
components of the graph built so far.

\writeup \emph{The rule.} Let $\sigma$ be a permutation and $a\ne b$. If $a$ and $b$ are in different cycles
$(a\,x_1\,\dots\,x_r)$ and $(b\,y_1\,\dots\,y_s)$, then $\sigma(a\,b)$ has the cycle $(a\,y_1\,\dots\,y_s\,b\,x_1\,\dots\,x_r)$ instead (check:
$a\mapsto b\mapsto y_1$, \dots, $b\mapsto a\mapsto x_1$): a merge. If they are in the same cycle, the same computation read backwards is a
split. So the number of cycles changes by $\pm1$.

(a) Let $\pi_i=\tau_1\cdots\tau_i$, $\pi_0=e$, with $n$ cycles. $\pi_{n-1}$ has one cycle if and only if every step merges. Let $T_i$ be
the graph with edges $\{a_1,b_1\},\dots,\{a_i,b_i\}$. If every step up to $i$ merges, then by induction the cycles of $\pi_i$ have the
same supports as the components of $T_i$: a merge joins the cycles of $a_{i+1}$ and $b_{i+1}$, and the new edge joins their
components. So step $i+1$ merges if and only if $a_{i+1}$ and $b_{i+1}$ are in different components of $T_i$, that is, if the new
edge creates no cycle. All $n-1$ steps merge if and only if $T_{n-1}$ is acyclic with $n-1$ edges, that is, a tree.

(b) By (a), the pairs (tree, ordering of its edges) are in bijection with the factorisations of $n$-cycles into $n-1$
transpositions: there are $n^{n-2}(n-1)!$ of them, by Cayley's formula. Conjugating by a permutation $\pi$ maps factorisations of
$c$ bijectively onto those of $\pi c\pi^{-1}$, and all $(n-1)!$ $n$-cycles are conjugate. So each has $\frac{n^{n-2}(n-1)!}{(n-1)!}=n^{n-2}$
factorisations.
\tools Volume~2, Chapter~4 (cycles, conjugacy); Volume~3, Chapter~12 (Cayley's formula).
\leads Dénes (1959). Counting factorisations of permutations into transpositions gives the \emph{Hurwitz numbers}, which count
branched covers of the sphere; Goulden and Jackson found a formula for every cycle type in 1997. Figure~\ref{fig:tree} shows one
factorisation.
\end{solution}

\begin{solution}{pb:conics}
\idea A square completion reduces every quadratic character sum to the one for $y^2-D$, and $y^2-z^2=D$ factors as
$(y-z)(y+z)=D$, which is easy to count.

\writeup (a) Since $p$ is odd, $4a\ne0$ and $\chi(4a)=\chi(a)$. From $4a(ax^2+bx+c)=(2ax+b)^2-D$ we get
$\chi(ax^2+bx+c)=\chi(a)\chi\bigl((2ax+b)^2-D\bigr)$, and $y=2ax+b$ runs over $\mathbb F_p$ once when $x$ does. So the sum is
$\chi(a)\sum_y\chi(y^2-D)$. If $D=0$ this is $\chi(a)(p-1)$. If $D\ne0$, the number of $z$ with $z^2=w$ is $1+\chi(w)$, so
\[
\sum_y\chi(y^2-D)=\#\{(y,z):y^2-z^2=D\}-p.
\]
With $u=y-z$, $v=y+z$ (an invertible change of variables, $2$ being invertible) the equation is $uv=D$, which has $p-1$
solutions. So the sum is $-1$, and the total is $-\chi(a)$.

(b) Suppose first $a\ne0$. For fixed $y$, the equation $ax^2+(by)x+(cy^2-1)=0$ in $x$ has $1+\chi(\Delta_y)$ solutions, where
$\Delta_y=b^2y^2-4a(cy^2-1)=Dy^2+4a$. The polynomial $Dy^2+4a$ has leading coefficient $D\ne0$ and discriminant $-16aD\ne0$, so
by (a)
\[
N=\sum_y\bigl(1+\chi(Dy^2+4a)\bigr)=p-\chi(D).
\]
If $a=0$, then $D=b^2\ne0$ and $\chi(D)=1$. The equation is $y(bx+cy)=1$: any $y\ne0$, and then $x=(y^{-1}-cy)/b$ is unique. That is
$p-1=p-\chi(D)$ solutions.

(c) For $n\ne0$, $Q=n$ is the equation $Q'=1$ with $Q'=Q/n$, whose discriminant $D/n^2$ has the same character. So $Q=n$ has
$p-\chi(D)\ge p-1>0$ solutions. For $x^2+y^2$, $D=-4$ and $\chi(-4)=\chi(-1)$.

\tools Volume~1, Chapter~6 (Legendre symbol).
\leads $x^2+y^2=1$ is a circle with $p+1$ or $p-1$ points: the $p-1$ case corresponds to a circle with two points at
infinity, $x\pm iy=0$, defined over $\mathbb F_p$. Every smooth conic has $p+1$ points in the projective plane, and the problem
is the affine shadow of that fact.
\end{solution}

\begin{solution}{pb:hashing}
\idea Two distinct points determine a line, so the values of a random line $ax+b$ at two distinct points are an arbitrary pair.
Reducing modulo $m$ loses a little uniformity, and that is all.

\writeup (a) For $x\ne y$ in $\Z/p\Z$ the linear map $(a,b)\mapsto(ax+b,ay+b)$ has determinant $x-y\ne0$, so it is a bijection of
$(\Z/p\Z)^2$, and maps the uniform law to the uniform law.

(b) By (a), $P(h(x)=h(y))$ is the proportion of pairs $(r,s)\in\{0,\dots,p-1\}^2$ with $r\equiv s\pmod m$. For each $r$ there are at most
$\lceil p/m\rceil\le\frac{p-1}m+1$ values of $s$ in its class, so the probability is at most $\frac1p\bigl(\frac{p-1}m+1\bigr)\le\frac1m+\frac1p$.

(c) The expected number of colliding pairs among $n$ keys is at most $\binom n2(\frac1{n^2}+\frac1p)\le\frac12+\frac{n^2}{2p}$. By Markov's
inequality the probability of at least one collision is at most that, so no collision happens with probability at least
$\frac12-\frac{n^2}{2p}$.

\tools Volume~0, Chapter~7 (linear maps); Volume~3, Chapters~2 and~4 (linearity of expectation, Markov's inequality).
\leads (c) is the first level of the Fredman--Komlós--Szemerédi perfect hashing scheme (1984), which stores $n$ keys in $O(n)$ space
with worst-case constant lookup time. Pairwise independence is enough here, which is why two random numbers $a,b$ suffice instead
of a fully random function.
\end{solution}

\begin{solution}{pb:triangleseigen}
\idea Closed walks of length $3$ are triangles traversed in one of six ways, and the eigenvalues add up their counts. The third
moment of the eigenvalues is bounded by the largest eigenvalue times the second moment.

\writeup (a) $A$ is real symmetric, so $\operatorname{tr}A^k=\sum_i\lambda_i^k$. Now $(A^2)_{vv}=d_v$, so $\operatorname{tr}A^2=\sum_vd_v=2m$; and $(A^3)_{vv}$ counts closed
walks $v\to x\to y\to v$, that is ordered pairs $(x,y)$ such that $vxy$ is a triangle, so every triangle is counted $6$ times in
$\operatorname{tr}A^3$.

(b) $6t=\sum\lambda_i^3\le\sum\lambda_i^2\,|\lambda_i|\le\lambda_{\max}\sum\lambda_i^2$, where $\lambda_{\max}=\max|\lambda_i|$, and
$\lambda_{\max}^2\le\sum\lambda_i^2=2m$. So $6t\le(2m)^{3/2}$, that is $t\le\frac{\sqrt2}3m^{3/2}$. For $K_q$, $t=\binom q3\sim\frac{q^3}6$ and
$m=\binom q2\sim\frac{q^2}2$, so $\frac t{m^{3/2}}\to\frac{1/6}{(1/2)^{3/2}}=\frac{\sqrt2}3$.

\tools Volume~11, Appendix~A (adjacency eigenvalues); Volume~0, Chapter~7 (spectral theorem).
\leads The exact maximum number of triangles with $m$ edges is given by the Kruskal--Katona theorem, and it is attained by
taking edges as ``close to a clique'' as possible. Problem~\ref{pb:shearer}(d) proves the same bound with entropy, and
Problem~\ref{pb:goodman} gives a lower bound for two colours.
\end{solution}

\begin{solution}{pb:walks}
\idea Walks are moments of the spectral measure seen from the uniform vector, and convexity compares a moment with a power of the
first moment. For closed $4$-walks, the top eigenvalue alone suffices.

\writeup (a) Let $(u_i)$ be an orthonormal basis of eigenvectors, $A u_i=\lambda_iu_i$, and $\mathbf1=\sum_ic_iu_i$. Then $\sum c_i^2=n$,
$\mathbf1^{\mathsf T}A\mathbf1=\sum c_i^2\lambda_i=2m$, and the number of walks of length $k$ is $w_k=\mathbf1^{\mathsf T}A^k\mathbf1=\sum c_i^2\lambda_i^k$. For even
$k$, $x\mapsto x^k$ is convex on $\R$, and Jensen with the weights $\frac{c_i^2}n$ gives $\frac{w_k}n\ge\bigl(\frac{2m}n\bigr)^k$.

(b) By Rayleigh's principle (Volume~11, Lemma~A.1), $\lambda_1\ge\frac{\mathbf1^{\mathsf T}A\mathbf1}{\mathbf1^{\mathsf T}\mathbf1}=\bar d$. So
$\operatorname{tr}A^4=\sum\lambda_i^4\ge\lambda_1^4\ge\bar d^{\,4}$.

(c) In $P_p$, $\bar d=d=\frac{p-1}2$, and $(A^2)_{xy}$ is $d$ if $x=y$, $\frac{p-5}4$ if $x,y$ are adjacent and $\frac{p-1}4$ otherwise
(Problem~\ref{pb:paley}(a)). So
\[
\operatorname{tr}A^4=\sum_{x,y}(A^2)_{xy}^2=pd^2+pd\Bigl(\frac{p-5}4\Bigr)^2+pd\Bigl(\frac{p-1}4\Bigr)^2.
\]
Here $pd^2=O(p^3)$, and each of the last two terms is $\frac p2\cdot p\cdot\frac{p^2}{16}+O(p^3)=\frac{p^4}{32}+O(p^3)$, so together they
give $\frac{p^4}{16}+O(p^3)$. Since $d^4=\frac{p^4}{16}+O(p^3)$ as well, $\operatorname{tr}A^4=d^4+O(p^3)=d^4(1+O(\frac1p))$.

\tools Volume~0, Chapter~7 (spectral theorem), Theorem~4.2 (Jensen); Volume~11, Lemma~A.1; Problem~\ref{pb:paley}.
\leads (a) holds for odd $k$ too (Blakley and Roy, 1965), but the proof is harder. (b) is the case of the $4$-cycle in Sidorenko's
conjecture, still open in general. Chung, Graham and Wilson (1989) showed that a sequence of graphs of density $\frac12$ with
$\operatorname{tr}A^4=(\frac n2)^4(1+o(1))$ behaves like a random graph in many other ways (``quasi-randomness''); the Paley graphs are the
standard example.
\end{solution}

\begin{solution}{pb:littlewoodofford}
\idea Make all coefficients positive. Then raising one more sign from $-1$ to $+1$ increases the sum by at least $2$, so the sets of
indices with sum in a window of length $2$ form an antichain, and Sperner applies.

\writeup (a) Replacing $a_i$ by $-a_i$ and $\varepsilon_i$ by $-\varepsilon_i$ does not change the law of the sum, so we may assume $a_i\ge1$. For $S\subseteq[n]$ let
$\sigma(S)=\sum_{i\in S}a_i-\sum_{i\notin S}a_i$, the sum when $\varepsilon_i=1$ exactly for $i\in S$. If $S\subsetneq T$, then $\sigma(T)-\sigma(S)=2\sum_{i\in T\setminus S}a_i\ge2$, so
$\sigma(S)$ and $\sigma(T)$ are not both in the open interval $J$. Thus $\{S:\sigma(S)\in J\}$ is an antichain, of size at most $\binom n{\lfloor n/2\rfloor}$
by Sperner's theorem, and each $S$ has probability $2^{-n}$.

(b) With all $a_i=1$, $\sum\varepsilon_ia_i=2|S|-n$, and the interval
$J=]\,2\lfloor\frac n2\rfloor-n-1,\,2\lfloor\frac n2\rfloor-n+1[$ contains exactly the value $2\lfloor\frac n2\rfloor-n$, taken for the $\binom n{\lfloor n/2\rfloor}$
sets with $|S|=\lfloor\frac n2\rfloor$. By Stirling's formula, $\binom{2k}k2^{-2k}\sim\frac1{\sqrt{\pi k}}$, which is $\sqrt{\frac2{\pi n}}$ for $n=2k$, and the odd
case is the same since $\binom{2k+1}k2^{-2k-1}=\binom{2k}k2^{-2k}\cdot\frac{2k+1}{2k+2}$.

\tools Volume~3, Theorem~A.16 (LYM and Sperner); Volume~0, Exercise~5.7 (Stirling's formula).
\leads Littlewood and Offord met this question in 1943 while counting real roots of random polynomials. For coefficients that are
also distinct, Erdős and Moser conjectured, and Stanley proved with the hard Lefschetz theorem, a bound of order $n^{-3/2}$.
Tao and Vu's ``inverse'' theorems explain which $a_i$ give large concentration, and they are central in random matrix theory.
\end{solution}

\begin{solution}{pb:cubeiso}
\idea Split the cube along one coordinate. Edges inside $S$ are edges inside each half plus a matching between the halves, at most
the size of the smaller half; an entropy-like inequality closes the induction.

\writeup Write $\log=\log_2$ and $0\log0=0$. Induction on $n$; for $n=0$, $e(S)=0$. Let $S_0,S_1$ be the parts of $S$ with last coordinate $0$
and $1$, seen in $\{0,1\}^{n-1}$, with $|S_0|=a\ge|S_1|=b$ (or the reverse). Edges inside $S$ are edges inside $S_0$, inside $S_1$, and edges in
the last direction, which join $x0$ to $x1$ and number at most $b$. So, by induction, $e(S)\le\frac12(a\log a+b\log b)+b$, and it suffices
that $a\log a+b\log b+2b\le(a+b)\log(a+b)$. For $b=0$ this is clear; otherwise put $t=\frac ba\in]0,1]$ and divide by $a$: it becomes
$g(t)=(1+t)\log(1+t)-t\log t-2t\ge0$. Now $g(0^+)=0$, $g(1)=2-2=0$, and $g''(t)=\frac1{\ln2}\bigl(\frac1{1+t}-\frac1t\bigr)<0$: $g$ is concave on $[0,1]$
and vanishes at the ends, so $g\ge0$.

A subcube of dimension $k$ has $2^k$ vertices and $k2^{k-1}=\frac12\cdot2^k\cdot k$ edges.

\tools Volume~0, Chapter~4 (convexity).
\leads This edge-isoperimetric inequality (Harper 1964, Bernstein 1967, Hart 1976) says equivalently that a set of $2^k$ vertices has at least
$2^k(n-k)$ edges leaving it. It is the combinatorial shadow of the logarithmic Sobolev inequality on the cube, and it underlies
the analysis of Boolean functions (Kahn--Kalai--Linial).
\end{solution}

\begin{solution}{pb:erdosselfridge}
\idea A potential function measures how dangerous the position is for Breaker: the sum over live sets of $2^{-(\text{points Maker still
needs})}$. A completed set has potential $1$. Maker's move can only add what Breaker's greedy move just removed.

\writeup (a) At any moment, call a set $H\in\mathcal H$ \emph{alive} if Breaker has no point of $H$, give it weight $w(H)=2^{-u(H)}$ with $u(H)$ the number of
points of $H$ not claimed by Maker, and let $\Phi=\sum_{H\text{ alive}}w(H)$. Initially $\Phi=\sum2^{-|H|}<\frac12$. Maker's first move at a point $x$ doubles the
weights of the alive sets containing $x$, so afterwards $\Phi<1$.

Consider a round in which Breaker claims $y$ and then Maker claims $x$. Breaker chooses $y$ maximising $W(y)=\sum_{H\text{ alive},\,H\ni y}w(H)$ among
the free points. His move kills those sets: $\Phi$ decreases by $W(y)$. Maker's move then doubles the weights of the alive sets containing
$x$ (none contains $y$ now), increasing $\Phi$ by at most $W(x)$, computed before Breaker's move, and $W(x)\le W(y)$ by the choice of $y$. So $\Phi$
does not increase over a round, and stays $<1$ until the end. If Maker ever claimed all of some $H$, that set would be alive with
weight $2^0=1\le\Phi$. So Maker never completes a set.

(b) Take $V$ the edge set of $K_N$ and $\mathcal H$ the edge sets of the $\binom Nq$ copies of $K_q$, each of size $\binom q2$. If $N\le2^{q/2}$, then
$\binom Nq2^{-\binom q2}\le\frac{N^q}{q!}2^{-q(q-1)/2}\le\frac{2^{q/2}}{q!}<\frac12$ for $q\ge3$.

\tools Volume~3, Theorem~12.2 (the probabilistic method, whose ``derandomisation'' this is); Volume~0, Chapter~8.
\leads The potential is the conditional expectation, under random play from now on, of the number of sets Maker would complete,
and the proof is the method of conditional expectations. In the clique game the Erd\H os--Selfridge bound is sharp in the exponent: Beck (2002) showed that the largest clique Maker can
build on $K_N$ has $2\log_2N-2\log_2\log_2N+O(1)$ vertices, so Maker wins the $K_q$ game once $N\ge2^{q/2+o(q)}$.
\end{solution}

\begin{solution}{pb:shearer}
\idea Entropy is additive along a chain and decreases when we condition on more. In Shearer's inequality each coordinate is
counted $k$ times on the right, each time with less conditioning than on the left.

\writeup (a) Let $p(x)=P(X=x)$ and $p(y\mid x)=P(Y=y\mid X=x)$. Then $\log\frac1{p(x)p(y|x)}=\log\frac1{p(x)}+\log\frac1{p(y|x)}$; averaging over $(X,Y)$ gives
$H(X,Y)=H(X)+H(Y\mid X)$. With subadditivity this gives $H(Y\mid X)=H(X,Y)-H(X)\le H(Y)$: conditioning on one variable never
raises entropy. Apply this to the conditional law of $(Z,Y)$ given $X=x$: it gives $H(Y\mid X=x,Z)\le H(Y\mid X=x)$, and averaging over
$x$ with the weights $P(X=x)$ gives $H(Y\mid X,Z)\le H(Y\mid X)$. Finally, Gibbs's inequality
against the uniform law gives $H(X)\le\log_2$(number of values).

(b) Iterating (a), $H(X)=\sum_iH(X_i\mid X_1,\dots,X_{i-1})$ (chain rule), and similarly, listing $F_j$ in increasing order,
\[
H(X_{F_j})=\sum_{i\in F_j}H\bigl(X_i\mid X_{i'},\ i'\in F_j,\ i'<i\bigr)\ge\sum_{i\in F_j}H(X_i\mid X_1,\dots,X_{i-1}),
\]
by (a), since $\{i'\in F_j:i'<i\}\subseteq\{1,\dots,i-1\}$. Summing over $j$, each $i$ appears at least $k$ times and the terms are non-negative, so
$\sum_jH(X_{F_j})\ge k\sum_iH(X_i\mid X_{<i})=kH(X)$.

(c) Let $X=(X_1,X_2,X_3)$ be uniform on $S$, so $H(X)=\log|S|$. The sets $\{1,2\},\{2,3\},\{1,3\}$ cover each coordinate twice, and $X_{\{1,2\}}$ takes
values in $S_{xy}$, so $H(X_{\{1,2\}})\le\log|S_{xy}|$, and similarly. By (b), $2\log|S|\le\log|S_{xy}|+\log|S_{yz}|+\log|S_{xz}|$.

(d) Let $(X_1,X_2,X_3)$ be uniform among the $6t$ ordered triples of vertices forming a triangle. Each pair $(X_i,X_j)$ is an ordered edge, one of
$2m$ values. By (b) with the same three pairs, $2\log(6t)\le3\log(2m)$, that is $6t\le(2m)^{3/2}$.

\tools Volume~3, Exercise~4.9 (subadditivity), Proposition~4.16 (Gibbs's inequality) and Chapter~4 (entropy).
\leads Shearer's lemma is due to Chung, Frankl, Graham and Shearer (1986). The same method proves Brègman's bound on permanents
(Radhakrishnan, 1997) and Kahn's bound on independent sets in regular bipartite graphs. Loomis and Whitney proved their inequality
for volumes of bodies in $\R^n$; it implies the isoperimetric inequality up to a constant.
\end{solution}

\begin{solution}{pb:sincpower}
\idea $\frac{\sin t}t=\mathbb E[e^{itV}]$ for $V$ uniform on $[-1,1]$. So $(\frac{\sin t}t)^n$ is the characteristic function of a sum of $n$ such
variables, and the inversion formula at $0$ turns the integral into a value of the density, which Problem~\ref{pb:irwinhall} computes.

\writeup (a) $\mathbb E[e^{itV}]=\frac12\int_{-1}^1e^{itv}\dd v=\frac{\sin t}t$. Let $Y=V_1+\cdots+V_n$; $\varphi_Y(t)=(\frac{\sin t}t)^n$ is integrable for $n\ge2$, so by
Volume~3, Theorem~6.8(b), $Y$ has a continuous density $f$ with $f(0)=\frac1{2\pi}\int(\frac{\sin t}t)^n\dd t$. Now $V_i=2U_i-1$ with $U_i$ uniform on $[0,1]$,
so $Y=2S_n-n$ and $f(0)=\frac12f_{S_n}(\frac n2)$, where $f_{S_n}(x)=\frac1{(n-1)!}\sum_j(-1)^j\binom nj(x-j)_+^{n-1}$ is the derivative of Problem~\ref{pb:irwinhall}(a). So
\[
\int\Bigl(\frac{\sin t}t\Bigr)^n\dd t=\pi f_{S_n}\Bigl(\frac n2\Bigr)=\frac\pi{(n-1)!}\sum_{j<n/2}(-1)^j\binom nj\Bigl(\frac n2-j\Bigr)^{n-1},
\]
which is the formula. For $n=2,3,4$: $\frac\pi{2}\cdot2=\pi$, $\frac\pi{8}(9-3)=\frac{3\pi}4$, $\frac\pi{48}(64-32)=\frac{2\pi}3$.

(b) On $[-\pi,\pi]$, $\frac{\sin t}t>0$ except at $\pm\pi$, and $\phi(t)=\log\frac{\sin t}t$ is smooth on $]-\pi,\pi[$ with a unique maximum $\phi(0)=0$ and
$\phi(t)=-\frac{t^2}6+O(t^4)$. By Laplace's method (Volume~4, Theorem~11.9), $\int_{-\pi}^\pi e^{n\phi}\sim\sqrt{\frac{2\pi}{n/3}}=\sqrt{\frac{6\pi}n}$. Outside,
$\bigl|\int_{|t|>\pi}(\frac{\sin t}t)^n\bigr|\le2\int_\pi^\infty t^{-n}\dd t=\frac{2\pi^{1-n}}{n-1}$, which is negligible.

\tools{Volume~3, Theorem~6.8 (inversion); Volume~4, Theorem~11.9 (Laplace's method); Problem~\ref{pb:irwinhall}.}
\leads (b) says $f(0)\sim\frac1{\sqrt{2\pi n/3}}$, the value at $0$ of the normal density of variance $\frac n3$: a local central limit theorem.
Integrals of this kind measure sections of the cube; with them Hensley (1979) and Ball (1986) proved that every central hyperplane
section of the unit cube of $\R^n$ has volume between $1$ and $\sqrt2$.
\end{solution}

\begin{solution}{pb:cantor}
\idea A sum of independent digits has a product characteristic function. Multiplying $t$ by $3$ shifts the product by one factor,
so along $t=2\pi3^m$ the function repeats itself and cannot decay; a density would force decay.

\writeup (a) Let $X_K=\sum_{k\le K}2B_k3^{-k}$. By independence,
\[
\mathbb E[e^{itX_K}]=\prod_{k\le K}\frac{1+e^{2it/3^k}}2=\prod_{k\le K}e^{it/3^k}\cos\frac t{3^k}=e^{it(1-3^{-K})/2}\prod_{k\le K}\cos\frac t{3^k}.
\]
As $|X-X_K|\le3^{-K}$, $|e^{itX}-e^{itX_K}|\le|t|3^{-K}$, so the left side tends to $\varphi(t)$; hence the product converges and the formula holds.

(b) For $t=2\pi3^m$, the factors with $k\le m$ are $\cos(2\pi3^{m-k})=1$, and the others are $\cos(2\pi3^{-j})$ with $j=k-m\ge1$. So
$|\varphi(2\pi3^m)|=\bigl|\prod_{j\ge1}\cos\frac{2\pi}{3^j}\bigr|=|\varphi(2\pi)|$. The first factor is $-\frac12$. For $j\ge2$, $x_j=\frac{2\pi}{3^j}\le\frac{2\pi}9<\frac\pi4$, and for
$|x|\le\frac\pi4$, $\cos x\ge1-\frac{x^2}2\ge e^{-x^2}$ (as $1-u\ge e^{-2u}$ for $0\le u\le\frac12$). So $\prod_{j\ge2}\cos x_j\ge\exp(-\sum_{j\ge2}x_j^2)=e^{-\pi^2/18}>0$,
and $|\varphi(2\pi)|>0$; numerically it is $0.3714\ldots$

If $X$ had a density $f$, then $\varphi(t)=\int f(x)e^{itx}\dd x\to0$ as $t\to\infty$ (Riemann--Lebesgue): for $f=1_{[a,b]}$, $|\int_a^be^{itx}\dd x|\le\frac2{|t|}$; this
extends to step functions by linearity, and to every integrable $f$ because step functions approximate $f$ in mean,
$\int|f-s|\le\varepsilon$, and $|\int(f-s)e^{itx}|\le\varepsilon$. But $\varphi(2\pi3^m)$ does not tend to $0$. So $X$ has no density.

(c) From $\sin2u=2\sin u\cos u$, by induction
\[
\sin t=2^m\sin\frac t{2^m}\prod_{k=1}^m\cos\frac t{2^k},
\]
and $2^m\sin\frac t{2^m}\to t$. So $\prod_{k\ge1}\cos\frac t{2^k}=\frac{\sin t}t$ for $t\ne0$. As in (a), $Y=\sum_k\varepsilon_k2^{-k}$ has characteristic function
\[
\prod_k\cos\frac t{2^k}=\frac{\sin t}t=\frac12\int_{-1}^1e^{itu}\dd u,
\]
that of the
uniform law on $[-1,1]$; by uniqueness (Volume~3, Theorem~6.8), $Y$ is uniform.

\tools{Volume~3, Exercise~5.7 and Theorem~6.8 (the Cantor law, characteristic functions); Volume~0, Chapter~2 (trigonometry).}
\leads The same question for $Y_\lambda=\sum_k\varepsilon_k\lambda^k$, the \emph{Bernoulli convolutions}, is a famous problem. For $\lambda<\frac12$ the law lives on a
Cantor set; for $\lambda=\frac12$ it is uniform. For $\frac12<\lambda<1$, Erdős (1939) used exactly the argument of (b) to prove that the law is
singular when $\frac1\lambda$ is a Pisot number, such as the golden ratio: then $\theta^m$ is very close to an integer, $\theta=\frac1\lambda$, so the
characteristic function does not decay along $t=2\pi\theta^m$. Solomyak (1995) proved that the law has a density for almost every
$\lambda\in\left]\frac12,1\right[$; which $\lambda$ exactly is still open.
\end{solution}

\begin{solution}{pb:majoritystable}
\idea Noise stability weights the spectrum by $\rho^{|S|}$, so it is a generating function of the spectral distribution. For majority,
the central limit theorem turns the two vote counts into a Gaussian pair, and the probability that two correlated Gaussians have
the same sign is an angle.

\writeup (a) $\mathbb E[x^Sy^T]=\prod_i\mathbb E[x_i^{1_S(i)}y_i^{1_T(i)}]$ by independence of the pairs; a factor with $i$ in exactly one of $S$, $T$ is $0$, and a
factor with $i\in S\cap T$ is $\rho$. So $\mathbb E[x^Sy^T]=\rho^{|S|}1_{S=T}$, and expanding $f(x)f(y)$ gives $\operatorname{Stab}_\rho(f)=\sum_S\rho^{|S|}\hat f(S)^2$.

(b) Let $X$, $Z$ be independent $N(0,1)$ and $Y=\rho X+\sqrt{1-\rho^2}Z$; then $(X,Y)$ is a Gaussian pair with the right covariance, so it has the right
law (Volume~3, Theorem~6.15). For a rotation $A$ of the plane, $A(X,Z)$ is Gaussian with covariance $AA^{\mathsf T}=I$, so it has the same law as
$(X,Z)$. Rotations permute the open sectors with vertex $0$ and a given angle, so the probability of such a sector depends only on its
angle $\theta$; it is additive in $\theta$, and the full plane has probability $1$, so it equals $\frac\theta{2\pi}$ for rational multiples of $2\pi$, and by
monotonicity for all $\theta$. Write $\rho=\cos\beta$ with $0<\beta<\pi$ and $(X,Z)=R(\cos\Theta,\sin\Theta)$; then $Y=R\cos(\Theta-\beta)$, and $\{X>0,Y>0\}$ is the sector
$\{|\Theta|<\frac\pi2,\ |\Theta-\beta|<\frac\pi2\}$ of angle $\pi-\beta$. So
\[
P(X>0,Y>0)=\frac{\pi-\arccos\rho}{2\pi}=\frac14+\frac{\arcsin\rho}{2\pi},
\]
and $\mathbb E[\operatorname{sgn}X\operatorname{sgn}Y]=2P(\operatorname{sgn}X=\operatorname{sgn}Y)-1=4P(X>0,Y>0)-1=\frac2\pi\arcsin\rho$.

(c) The vectors $(x_i,y_i)$ are independent and identically distributed, with mean $0$ and covariance matrix $\begin{pmatrix}1&\rho\\\rho&1\end{pmatrix}$, so
$\frac1{\sqrt n}\sum_i(x_i,y_i)$ converges in law to $(X,Y)$. For $n$ odd the sums are never $0$, and $\operatorname{Stab}_\rho(\operatorname{Maj}_n)=4P(\sum x_i>0,\sum y_i>0)-1$ by the
symmetry $(x,y)\mapsto(-x,-y)$. The quadrant $\{u>0,v>0\}$ has a boundary of probability $0$ for $(X,Y)$, so its probability converges, and
$\operatorname{Stab}_\rho(\operatorname{Maj}_n)\to\frac2\pi\arcsin\rho$.

(d) Let $G_n(\rho)=\operatorname{Stab}_\rho(\operatorname{Maj}_n)=\sum_kW^k_n\rho^k$, where $W^k_n=W^k(\operatorname{Maj}_n)\ge0$ and $\sum_kW^k_n=1$, and $G(\rho)=\frac2\pi\arcsin\rho=\sum_kb_k\rho^k$ with
$b_{2j+1}=\frac2\pi\binom{2j}j\frac1{4^j(2j+1)}$ and $b_{2j}=0$ (integrate $(1-u^2)^{-1/2}=\sum_j\binom{2j}j\frac{u^{2j}}{4^j}$). Take any subsequence; by a diagonal extraction, a
further subsequence has $W^k_n\to w_k$ for every $k$. For $0\le\rho<1$, the tail $\sum_{k\ge K}W^k_n\rho^k$ is at most $\frac{\rho^K}{1-\rho}$, uniformly in $n$, so
$\sum_kw_k\rho^k=G(\rho)$ on $[0,1[$. Two power series that agree on $[0,1[$ have the same coefficients (let $\rho\to0$, divide by $\rho$, repeat), so
$w_k=b_k$. Every subsequence has a further subsequence with these limits, so $W^k_n\to b_k$. In particular $W^1\to b_1=\frac2\pi$, as in
Problem~\ref{pb:influence}, and $W^3\to\frac1{3\pi}$.

\tools{Volume~3, Theorem~6.15 and the remark after it (Gaussian vectors, the CLT in $\R^2$); Volume~4, Chapter~3 (power series).}
\leads Sheppard found his formula in 1899, studying correlations in biometric data. The limit $\frac2\pi\arcsin\rho$ is the starting point of
the \emph{Majority is Stablest} theorem of Mossel, O'Donnell and Oleszkiewicz (2010): among Boolean functions with mean $0$ and all
influences small, none is more noise-stable than majority. With the unique games conjecture, it implies that the Goemans--Williamson
algorithm for maximum cut, whose approximation ratio $0.878\ldots$ comes from the same arcsine, is optimal.
\end{solution}

\begin{solution}{pb:period}
\idea Modulo a normal subgroup containing all the ratios of steps, the walk moves like a clock hand. Conversely the cyclic classes of
a periodic walk are the cosets of a normal subgroup.

\writeup (a) Let $\Sigma$ be the set of products of one or more elements of $S$. It is closed under products. In a finite group each $s$ has finite
order $o(s)$, so $e=s^{o(s)}$ and $s^{-1}=s^{o(s)-1}$ lie in $\Sigma\cup\{e\}$. So $\Sigma\cup\{e\}$ is a subgroup containing $S$, hence equal to $G$. Every $g$ is reachable from $e$,
and from $x$ the walk reaches $xg$ along the same steps: the walk is irreducible.

(b) Every step lies in $gN$, and $N$ is normal, so $X_t=\xi_1\cdots\xi_t\in g^tN$. The images of $S$ in $G/N$ all equal $gN$ and generate $G/N$, so $G/N$ is
cyclic, generated by $gN$, of order $m=[G:N]\ge2$. The walk can be at $e$ at time $t$ only if $g^t\in N$, that is $m\mid t$: the period is a multiple of $m$. In
particular $P(X_t\in N)=0$ when $m\nmid t$, while the uniform law gives $N$ the mass $\frac1m$; the walk does not converge.

(c) Let $N_r$ be the set of $x$ that the walk reaches at some time $t\equiv r\pmod d$. Since $e$ is reachable from every $x$, say in $u$ steps, two times $t,t'$
at which $x$ is reached satisfy $t+u,t'+u\in d\Z$; so the sets $N_0,\dots,N_{d-1}$ are disjoint, and they cover $G$ by (a). If $x\in N_0$ is reached at time $t$ and $y\in N_0$
at time $t'$, then $xy$ is reached at time $t+t'$ (go to $x$, then copy the path to $y$); so $N=N_0$ is closed under products, and being finite and
containing $e$ it is a subgroup.

Fix $s\in S$. The walk goes from $s$ back to $e$ in some $u$ steps, so $s^{-1}$ is reached at time $u$ with $1+u\equiv0$. If $x\in N$ is reached at time $t$, then $s^{-1}xs$
is reached at time $u+t+1\equiv0$: so $s^{-1}Ns\subseteq N$, hence $g^{-1}Ng\subseteq N$ for every $g$, since $S$ generates $G$ as a semigroup; $N$ is normal. If $x\in N_1$ is
reached at time $t\equiv1$, then $s^{-1}x$ is reached at time $u+t\equiv0$, so $x\in sN$; conversely $sN\subseteq N_1$. Thus $S\subseteq N_1=sN$, and likewise $N_r=s^rN$, so the $d$
classes are the cosets of $N$ and $[G:N]=d$.

If $S$ lies in no coset of a proper normal subgroup, then $N=G$ and $d=1$: the walk is irreducible and aperiodic, the uniform law is
stationary, and Volume~3, Theorem~8.11 gives convergence. If $S$ does lie in such a coset, (b) forbids convergence.

(d) All transpositions are odd, so $S$ lies in the coset $\tau A_n$ of the normal subgroup $A_n$ of index $2$. By (b) the period is even; since $\tau\tau=e$ the walk can
return at time $2$, so the period is $2$. The deck never converges: its sign alternates.

(e) The group $(\Z/p\Z)^*$ is abelian, so every subgroup is normal, and $\{2,3\}\subseteq gN$ means $2N=3N$, that is $q=3\cdot2^{-1}\in N$. The smallest such subgroup is $H=\langle q\rangle$. So
$S$ lies in a coset of a proper subgroup if and only if $q$ is not a primitive root, and then, by (b) and (c), the subgroup $N$ of (c) contains $H$
while the period is a multiple of $[G:H]$, so the period is exactly $[G:H]$. For $p=5$: $2^{-1}=3$ and $q=9\equiv4$, of order $2$ in a group of order $4$: period $2$
(indeed $2$ and $3$ are both non-squares modulo $5$, and the walk alternates between squares and non-squares). For $p=13$: $2^{-1}=7$, $q=21\equiv8$, and
$8^2\equiv12$, $8^4\equiv1$: $q$ has order $4$ in a group of order $12$, so the period is $3$. In both cases $2$ is a primitive root, so the hypothesis holds.

\tools{Volume~2, Chapters~1--2 (cosets, normal subgroups, quotients); Volume~3, Lemma~8.4 and Theorem~8.11.}
\leads The criterion is classical; Diaconis's book \emph{Group Representations in Probability and Statistics} (1988) states it in this form. How often
$3\cdot2^{-1}$ is a primitive root is a case of Artin's conjecture on primitive roots: Artin predicted in 1927 that $2$, for instance, is a
primitive root for about $37.4\%$ of all primes, and Hooley proved it in 1967 assuming the generalised Riemann hypothesis.
\end{solution}

\begin{solution}{pb:rollingdie}
\idea Tipping is a quarter turn, which is an odd permutation of the four diagonals, and a closed tour on a chessboard has even length. So
closed tours give even permutations only, and two small loops already give all of them.

\writeup Put the $x$-axis east, the $y$-axis north and the $z$-axis up, and describe the orientation of the cube by the rotation $R$ carrying
it from its initial position. A move replaces $R$ by $QR$, where $Q$ is one of the quarter turns
\[
R_E=\begin{pmatrix}0&0&1\\0&1&0\\-1&0&0\end{pmatrix},\quad R_N=\begin{pmatrix}1&0&0\\0&0&1\\0&-1&0\end{pmatrix},\quad R_W=R_E^{\mathsf T},\quad R_S=R_N^{\mathsf T}
\]
(rolling east sends the top face to the east, so $R_Ee_z=e_x$; similarly $R_Ne_z=e_y$).

(a) The rotation group of the cube permutes its four long diagonals, and this identifies it with $\mathfrak S_4$ (Volume~2, Proposition~1.5 and Exercise~1.6). A
quarter turn about the $z$-axis, $(x,y,z)\mapsto(-y,x,z)$, sends the diagonals through $(1,1,1)$, $(-1,1,1)$, $(-1,-1,1)$, $(1,-1,1)$ each to the next: it is a
$4$-cycle, an odd permutation (Volume~2, Corollary~4.3). Quarter turns about the other face axes are conjugate to it, hence also odd.

(b) Colour the board like a chessboard. Each move changes the colour of the square, so a tour that returns to its square has an even
number of moves, and its total rotation is a product of an even number of odd permutations: it is even. A quarter turn about the vertical
axis is a $4$-cycle, odd: impossible. A half turn is the square of a $4$-cycle, a product of two transpositions, even; by (c) it is possible.

(c) Multiplying the four matrices,
\[
R_SR_WR_NR_E=\begin{pmatrix}0&1&0\\0&0&-1\\-1&0&0\end{pmatrix},\qquad R_NR_WR_SR_E=\begin{pmatrix}0&-1&0\\0&0&1\\-1&0&0\end{pmatrix}.
\]
Both have trace $0=1+2\cos\theta$, so they are rotations by $\theta=\frac{2\pi}3$; their axes, the fixed vectors, are $(1,1,-1)$ and $(1,-1,-1)$: two different long diagonals.
As permutations of the diagonals they are $3$-cycles fixing different diagonals, so they generate a subgroup of $A_4$ with two different
subgroups of order $3$, of order divisible by $3$ and larger than $3$. As $A_4$ has no subgroup of order $6$ (Volume~2, Proposition~4.6), the subgroup is $A_4$
itself. Concatenating closed tours multiplies their rotations, so the rotations of closed tours form a group, which contains $A_4$ and is
contained in $A_4$ by (b).

(d) The orientation performs a random walk on the rotation group (with steps multiplied on the left, which changes nothing in
Problem~\ref{pb:period}). Its steps lie in the odd coset of the normal subgroup $A_4$, so the walk has period $2$ and does not converge. At even
times it is a walk on $A_4$ whose steps (two moves) include the identity ($R_WR_E=I$) and generate $A_4$ by (c); it converges to the uniform law on
$A_4$, and at odd times $X_{2t+1}$ converges to the uniform law on the odd coset. The body face on top is $R^{-1}e_z$. The group $A_4$ acts transitively on
the six faces $\pm e_x,\pm e_y,\pm e_z$: the rotations of order $3$ about $(1,1,1)$ permute $e_x,e_y,e_z$ cyclically, and the half turn about the $z$-axis sends $e_x$
to $-e_x$. By orbit--stabiliser (Volume~2, Theorem~3.5) every face is hit by exactly $2$ of the $12$ elements, so under the uniform law on $A_4$ (and also on
$A_4g$, which moves $g^{-1}e_z$ by a uniform element of $A_4$) each face is on top with probability $\frac16$. Both parities give $\frac16$, so the face on top converges.

(e) The orientation walk is irreducible on $24$ states (its steps and their products generate the whole group) and its stationary law is
uniform, so the expected return time is $24$ (Volume~3, Theorem~8.8).

\tools{Volume~2, Proposition~1.5, Exercise~1.6, Corollary~4.3, Proposition~4.6 and Theorem~3.5; Volume~3, Theorem~8.8; Problem~\ref{pb:period}.}
\leads The same parity count is the reason why half of the positions of the fifteen puzzle cannot be reached (Volume~2, Theorem~4.7).
The pair (square, orientation) is itself a random walk on a group, in which the translations of $\Z^2$ are twisted by the rotations of
the cube; a numerical calculation estimates how often the rolling die is back on its square in its initial orientation.
\end{solution}

\begin{solution}{pb:gapdiam}
\idea A function that varies a lot overall must vary along some edge: connect any two elements by a short path and pay for its
length twice, once in Cauchy--Schwarz and once in the number of paths through an edge.

\writeup (a) Expanding the square and using that $g\mapsto gs$ is a bijection,
$\sum_{g,s}(f(gs)-f(g))^2=2|S|\sum_gf(g)^2-2\sum_gf(g)\sum_sf(gs)$, so $\mathcal E(f)=\frac1{|G|}\sum_gf(g)\bigl(f(g)-(Pf)(g)\bigr)=\langle f,(I-P)f\rangle$. Similarly
$\operatorname{Var}(f)=\frac1{|G|}\sum f^2-(\frac1{|G|}\sum f)^2$. Since $S=S^{-1}$, $P$ is a symmetric matrix, hence self-adjoint for $\langle\cdot,\cdot\rangle$, with an orthonormal basis of eigenvectors, the
constant $\mathbf 1$ for $\lambda_1=1$. Neither $\mathcal E$ nor $\operatorname{Var}$ changes when a constant is added to $f$, so we may assume $f\perp\mathbf1$, and then $\frac{\mathcal E(f)}{\operatorname{Var}(f)}=\frac{\langle f,(I-P)f\rangle}{\langle f,f\rangle}\ge1-\lambda_2$,
with equality for an eigenvector of $\lambda_2$.

(b) For $x\in G$ fix a word $x=s_1\cdots s_\ell$ with $\ell\le d$, and put $w_i=s_1\cdots s_i$. Then $f(gx)-f(g)=\sum_{i=1}^\ell\bigl(f(gw_i)-f(gw_{i-1})\bigr)$, so by Cauchy--Schwarz
\[
\sum_g\bigl(f(gx)-f(g)\bigr)^2\le\ell\sum_{i=1}^\ell\sum_g\bigl(f(gw_{i-1}s_i)-f(gw_{i-1})\bigr)^2\le\ell^2\sum_{g'}\sum_{s\in S}\bigl(f(g's)-f(g')\bigr)^2,
\]
after the substitution $g'=gw_{i-1}$. The right side is at most $2d^2|G||S|\mathcal E(f)$. Summing over the $|G|$ values of $x$ (with $h=gx$),
$\operatorname{Var}(f)=\frac1{2|G|^2}\sum_x\sum_g(f(gx)-f(g))^2\le|S|d^2\mathcal E(f)$. With (a), $1-\lambda_2\ge\frac1{|S|d^2}$.

(c) For $\Z/n\Z$, $|S|=2$ and $d=\lfloor\frac n2\rfloor$, so $1-\lambda_2\ge\frac2{n^2}$, while the truth is $1-\cos\frac{2\pi}n\sim\frac{2\pi^2}{n^2}$: only the constant is lost. For $\mathfrak S_n$
with adjacent transpositions, $|S|=n-1$ and $d=\binom n2$ (an adjacent transposition changes the number of inversions by one, and bubble
sort shows that $\binom n2$ steps suffice), so $1-\lambda_2\ge\frac4{(n-1)^3n^2}$. The true gap is $\frac2{n-1}(1-\cos\frac\pi n)\sim\frac{\pi^2}{n^3}$ (a direct computation verifies it for $n\le6$): here the crude
path argument loses a factor of order $n^2$.

\tools{Volume~3, Proposition~8.16 (eigenvalues of reversible chains); spectral theorem for symmetric matrices (Volume~0); Cauchy--Schwarz.}
\leads Choosing the paths more cleverly, so that few of them use each edge, is the method of canonical paths of Jerrum and Sinclair and
the comparison theorems of Diaconis and Saloff-Coste (1993), which transfer a gap from a walk one understands to one that one
does not. For adjacent transpositions, Wilson (2004) showed that the mixing time is of order $n^3\ln n$.
\end{solution}

\begin{solution}{pb:treewalk}
\idea On the tree the distance to the root is a walk with a drift away from it, so returns are rare: the return probability decays
exponentially, at the rate $\frac{\sqrt3}2$, while on $\Z^2$ it decays only like $\frac1{\pi n}$.

\writeup (a) A reduced word of length $k\ge1$ ends with some letter $\ell$; multiplying by $\ell^{-1}$ shortens it, and by any of the three other letters
lengthens it, without cancellation. From $e$ every letter gives length $1$. The transition probabilities depend only on the length, so $(D_t)$ is
a Markov chain with the stated transitions.

(b) The walk returns to $e$ exactly when $D$ returns to $0$, which it does if and only if, from $1$, it ever reaches $0$. For the walk on $\Z$ with
up-probability $p=\frac34$ and down-probability $q=\frac14$ started at $1$, this happens with probability $\frac qp=\frac13$ (Volume~3, Corollary~7.10). The number of visits
to $e$ is geometric: $P(V\ge k)=3^{-(k-1)}$, so $\mathbb E[V]=\frac1{1-1/3}=\frac32$.

(c) For $k\ge1$, $\frac14h(k-1)+\frac34h(k+1)=h(k)\bigl(\frac{\sqrt3}4+\frac3{4\sqrt3}\bigr)=\frac{\sqrt3}2h(k)$, and for $k=0$, $h(1)=\frac1{\sqrt3}\le\frac{\sqrt3}2$. So $\mathbb E[h(D_{t+1})]\le\frac{\sqrt3}2\mathbb E[h(D_t)]$, and by induction
$\mathbb E[h(D_{2n})]\le(\frac{\sqrt3}2)^{2n}$. Since $h\ge0$ and $h(0)=1$, $P(D_{2n}=0)\le\mathbb E[h(D_{2n})]$.

(d) Consider the paths with $D_1=1$, $D_s\ge1$ for $1\le s\le2n-1$, $D_{2n-1}=1$ and $D_{2n}=0$. Between times $1$ and $2n-1$ they make $n-1$ up-steps and $n-1$ down-steps
without going below $1$, in $C_{n-1}$ ways ($C_m=\frac1{m+1}\binom{2m}m$, the Catalan number), each of probability $(\frac34)^{n-1}(\frac14)^{n-1}$, and the last step has probability
$\frac14$. With $\binom{2m}m\ge\frac{4^m}{2m+1}$ we get $C_{n-1}\ge\frac{4^{n-1}}{(2n-1)n}\ge\frac{4^{n-1}}{2n^2}$, hence
\[
P(D_{2n}=0)\ge\frac14\cdot\frac{4^{n-1}}{2n^2}\Bigl(\frac3{16}\Bigr)^{n-1}=\frac1{8n^2}\Bigl(\frac34\Bigr)^{n-1}.
\]
Taking $2n$-th roots, both bounds tend to
$\sqrt{3/4}=\frac{\sqrt3}2$.

(e) On the $q$-regular tree the distance goes up with probability $\frac{q-1}q$ and down with probability $\frac1q$ (from $k\ge1$). With $h(k)=(q-1)^{-k/2}$,
$\frac1qh(k-1)+\frac{q-1}qh(k+1)=\frac{2\sqrt{q-1}}qh(k)$, and $h(1)\le\frac{2\sqrt{q-1}}q$ since $q\ge2$; the path count gives $C_{n-1}(\frac{q-1}{q^2})^{n-1}\frac1q$, whose $2n$-th root tends to
$\frac{2\sqrt{q-1}}q$. On $\Z^2$, $P(S_{2n}=0)=u_{2n}^2\sim\frac1{\pi n}$ (two independent coordinates after a rotation by $45^\circ$, as in Problem~\ref{pb:coprimewalk}), whose $2n$-th
root tends to $1$.

\tools{Volume~3, Corollary~7.10 (gambler's ruin against an infinite bank) and Theorem~7.3; Catalan numbers (Volume~0).}
\leads Kesten proved in 1959 that the return probabilities of a symmetric walk on a group decay exponentially if and only if the group is
not \emph{amenable}, the property that fails for the free group and that underlies the Banach--Tarski paradox (Volume~7). The number
$2\sqrt{q-1}$, the spectral radius of the tree's adjacency operator, is the benchmark of the Ramanujan graphs of Volume~11, Appendix~A.
\end{solution}

\begin{solution}{pb:commute}
\idea A cycle that starts and ends at $a$ spends, on average, time at each vertex in proportion to the stationary law. Applied to the round
trip $a\to b\to a$, it equates two ways of counting visits to $a$: by the stationary law, and by the escape probability of the electrical
network.

\writeup (a) Let $\nu(x)=\mathbb E_a\bigl[\sum_{t<\tau}\mathbf 1_{X_t=x}\bigr]$, which is finite since $\nu(x)\le\mathbb E_a[\tau]$. For each $y$,
\[
\nu(y)=\mathbb E_a\Bigl[\sum_{t=1}^{\tau}\mathbf 1_{X_t=y}\Bigr]=\sum_{t\ge1}\sum_xP_a(X_{t-1}=x,\ \tau>t-1)P(x,y)=\sum_x\nu(x)P(x,y).
\]
The first equality holds because $X_0=X_\tau=a$, so shifting the window by one changes nothing; the second uses the Markov property at time $t-1$,
the event $\{\tau>t-1\}$ being determined by $X_0,\dots,X_{t-1}$. So $\nu$ is a non-negative invariant measure of total mass $\mathbb E_a[\tau]$. By uniqueness
(Volume~3, Theorem~8.8), $\nu=\mathbb E_a[\tau]\,\pi$.

(b) Here $\mathbb E_a[\tau]=\mathbb E_a[T_b]+\mathbb E_b[T_a]$ by the strong Markov property at $T_b$. The visits to $a$ before $\tau$ all happen before $T_b$: the walk starts at $a$,
makes excursions from $a$, and stops counting once an excursion reaches $b$. Each excursion reaches $b$ before returning to $a$ with
probability $\frac1{\deg(a)R(a,b)}$, independently (Volume~3, Exercise~7.8(b)), so the number of visits to $a$ is geometric with mean $\deg(a)R(a,b)$. By (a)
it equals $\pi(a)\mathbb E_a[\tau]=\frac{\deg a}{2|\mathcal E|}\mathbb E_a[\tau]$, and the identity follows.

(c) On a path $0,1,\dots,m$, $R=m$ (resistors in series) and $|\mathcal E|=m$, so the identity predicts a commute time $2m^2$. Indeed the walk on the path
started at $0$ is $|S|$ for the simple walk $S$ on $\Z$, so $T_m$ is the exit time of $S$ from $\{-m+1,\dots,m-1\}$, of mean $m\cdot m=m^2$ by the gambler's ruin formula
$\mathbb E_x[\tau]=x(N-x)$ (Volume~3, Theorem~7.9, with $N=2m$, $x=m$); by symmetry the trip back also takes $m^2$. On a cycle of length $n$, two neighbours are joined by
resistances $1$ and $n-1$ in parallel, so $R=\frac{n-1}n$ and $2|\mathcal E|R=2(n-1)$. Directly: unrolling the cycle, the walk from a neighbour of $b$ hits $b$ when it
leaves $\{1,\dots,n-1\}$ from $1$, which takes $1\cdot(n-1)$ on average; each one-way trip takes $n-1$.

(d) If $a\sim b$, the edge $ab$ alone carries a unit flow of energy $1$, so $R(a,b)\le1$ by Thomson's principle (Problem~\ref{pb:rayleigh}(a)). Then $\mathbb E_a[T_b]<\mathbb E_a[T_b]+\mathbb E_b[T_a]\le2|\mathcal E|$. Take a spanning tree and a closed walk
$v_0,v_1,\dots,v_{2n-2}=v_0$ round it that uses each of its $n-1$ edges once in each direction. Following the walk from $v_0$, wait each time until the
random walk hits the next vertex $v_{i+1}$; when the last one is reached, every vertex has been visited. The expected time is at most
$\sum_i\mathbb E_{v_i}[T_{v_{i+1}}]=\sum_{\text{tree edges}}(\mathbb E_x[T_y]+\mathbb E_y[T_x])\le(n-1)\cdot2|\mathcal E|$, where the tour is started at the walk's starting vertex $v_0$.

\tools{Volume~3, Theorem~8.8, Corollary~8.15, Exercise~7.8 and Theorem~7.9.}
\leads The commute time identity is due to Chandra, Raghavan, Ruzzo, Smolensky and Tiwari (1989), and the cover time bound to Aleliunas,
Karp, Lipton, Lovász and Rackoff (1979), who used it to decide connectivity in logarithmic space with a random walk. In 2005 Reingold
removed the randomness: undirected connectivity can be decided deterministically in logarithmic space.
\end{solution}

\begin{solution}{pb:foster}
\idea The resistances are the entries of the inverse of the Laplacian on the hyperplane $\mathbf 1^\perp$, and summing them over the edges rebuilds
the Laplacian inside a trace, whose value is the dimension $n-1$.

\writeup (a) $L$ is symmetric and $\langle x,Lx\rangle=\sum_{\{u,w\}\in\mathcal E}(x_u-x_w)^2$, which vanishes only for constant $x$ ($G$ is connected). So $\ker L=\R\mathbf 1$, and the image of $L$ is
$(\ker L)^\perp=W$. As $e_a-e_b\in W$, there is $v$ with $Lv=e_a-e_b$, unique up to constants. Let $h$ be the voltage of Volume~3, Exercise~7.8: $h(a)=0$,
$h(b)=1$, harmonic elsewhere, and $I=\sum_{y\sim a}h(y)=-(Lh)(a)$. Then $(Lh)(x)=0$ for $x\ne a,b$, and $\sum_x(Lh)(x)=\langle\mathbf1,Lh\rangle=0$ gives $(Lh)(b)=I$. So $v=-\frac hI$
satisfies $Lv=e_a-e_b$, and $v(a)-v(b)=\frac1I=R(a,b)$.

(b) Let $M$ be the inverse of $L$ restricted to $W$. For $x\sim y$, $R(x,y)=\langle e_x-e_y,M(e_x-e_y)\rangle$, since $v=M(e_x-e_y)$ is a solution and $\langle e_x-e_y,v\rangle=v(x)-v(y)$. Write
$\langle u,Mu\rangle=\operatorname{tr}(M\,uu^{\mathsf T})$ for $u\in W$, and note that $L=\sum_{\{x,y\}\in\mathcal E}(e_x-e_y)(e_x-e_y)^{\mathsf T}$, as operators on $W$. Then
\[
\sum_{\{x,y\}\in\mathcal E}R(x,y)=\operatorname{tr}_W\Bigl(M\sum_{\{x,y\}}(e_x-e_y)(e_x-e_y)^{\mathsf T}\Bigr)=\operatorname{tr}_W(ML)=\dim W=n-1.
\]

(c) An automorphism $\gamma$ maps the voltage problem for $(a,b)$ to the one for $(\gamma a,\gamma b)$, so $R(\gamma a,\gamma b)=R(a,b)$. If the automorphisms act transitively
on edges, all edges have the same resistance, and Foster's theorem gives $|\mathcal E|R=n-1$. If some automorphism exchanges $x$ and $y$, it
exchanges the laws of the walks from $x$ and from $y$, so $\mathbb E_x[T_y]=\mathbb E_y[T_x]$, and the commute time identity gives $2\mathbb E_x[T_y]=2|\mathcal E|R=2(n-1)$.

(d) The cube, the Petersen graph and the dodecahedron have the required symmetry (their automorphism groups act transitively on
ordered pairs of neighbours). So the expected times are $n-1$: $7$, $9$ and $19$.

\tools{Volume~3, Exercise~7.8; Problem~\ref{pb:commute}; Volume~0 (symmetric matrices, trace); Volume~2, Chapter~3 (group actions).}
\leads Foster published the theorem in 1949. It has a one-line probabilistic meaning: for an edge $xy$, $R(x,y)$ is the probability that the
edge belongs to a uniformly chosen spanning tree (a consequence of Kirchhoff's work of 1847), and every spanning tree has $n-1$ edges.
Problem~\ref{pb:aldousbroder} shows how a random walk samples such a tree.
\end{solution}

\begin{solution}{pb:rayleigh}
\idea Current is the unit flow of least energy. Deleting edges shrinks the set of flows, so it can only increase the minimum; and any
family of disjoint cuts forces every flow to pay a toll at each cut. In the plane the tolls add up to a divergent harmonic series.

\writeup (a) $i(x,y)=v(x)-v(y)$ is antisymmetric, and $\sum_yi(x,y)=\deg(x)v(x)-\sum_{y\sim x}v(y)=(Lv)(x)$, which is $1$ at $a$, $-1$ at $b$ and $0$ elsewhere: a unit flow. Its
energy is $\sum_{\{x,y\}}(v(x)-v(y))^2=\langle v,Lv\rangle=v(a)-v(b)=R(a,b)$. If $\theta$ is another unit flow, $\eta=\theta-i$ has zero net flow at every vertex, and
\[
\sum_{\{x,y\}}i(x,y)\eta(x,y)=\frac12\sum_x\sum_{y\sim x}(v(x)-v(y))\eta(x,y)=\sum_xv(x)\sum_{y\sim x}\eta(x,y)=0,
\]
using antisymmetry of $\eta$ for the middle equality. Hence $\mathcal W(\theta)=\mathcal W(i)+\mathcal W(\eta)\ge R(a,b)$.

(b) A unit flow on the graph with fewer edges is a unit flow on the larger graph (put $0$ on the removed edges), with the same energy. The
minimum over the smaller set is larger.

(c) Let $\theta$ be a unit flow and $A$ the set of vertices that $a$ can reach without using an edge of $\Pi_k$, so $b\notin A$. Summing the net flows over $A$,
the contributions of edges inside $A$ cancel, and $1=\sum_{x\in A,\,y\notin A,\,x\sim y}\theta(x,y)$; all those edges belong to $\Pi_k$. So $\sum_{e\in\Pi_k}|\theta_e|\ge1$, and by Cauchy--Schwarz
$\sum_{e\in\Pi_k}\theta_e^2\ge\frac1{|\Pi_k|}$. The $\Pi_k$ are disjoint, so $\mathcal W(\theta)\ge\sum_k\frac1{|\Pi_k|}$; minimise over $\theta$.

(d) For $0\le k<n$ let $\Pi_k$ be the set of edges from $\partial B_k$ to $\partial B_{k+1}$ (with $\partial B_0=\{0\}$). There are $4(2k+1)$ of them (each side of the square $\partial B_k$ has
$2k+1$ points with one outward edge), they are disjoint, and every path from $0$ to $z_n$ crosses each of them. By (c),
\[
R(0,z_n)\ge\sum_{k=0}^{n-1}\frac1{4(2k+1)}\ge\frac14\int_0^n\frac{\dd x}{2x+1}=\frac18\ln(2n+1).
\]

(e) Let $H_n$ be the graph with the edges of $H$ inside $B_n$, the boundary merged into $z_n$. If $0$ is not connected to $z_n$ in $H_n$, the component of $0$ in $H$
stays inside $B_n$; it is finite, the walk on it is a finite irreducible chain, and it returns to $0$ for sure (Volume~3, Theorem~8.8). Otherwise, working in
the component of $0$ in $H_n$ and using Rayleigh, $R_{H_n}(0,z_n)\ge R(0,z_n)\ge\frac18\ln(2n+1)$, and by Volume~3, Exercise~7.8(b), the walk on $H$
started at $0$ reaches $\partial B_n$ before returning to $0$ with probability $\frac1{\deg_H(0)R_{H_n}(0,z_n)}\le\frac8{\ln(2n+1)}$ (until it reaches $\partial B_n$ the two walks coincide). Now
$P(\text{no return})\le P(\text{reach }\partial B_n\text{ before returning})+P(\text{stay in }B_n\text{ for ever without returning})$. The second term is $0$: from every vertex of $B_n$ that the walk
can reach there is a path in $H$ to $0$, which stays in $B_n$ or touches $\partial B_n$, so within $|B_n|$ steps the walk hits $\{0\}\cup\partial B_n$ with probability at least some
$\delta>0$, and it does so eventually with probability $1$ (as in Volume~3, Lemma~7.8). Letting $n\to\infty$, $P(\text{no return})=0$.

\tools{Volume~3, Exercise~7.8, Lemma~7.8, Theorem~8.8; Problem~\ref{pb:foster}(a); Cauchy--Schwarz.}
\leads The Nash-Williams criterion (1959) and Thomson's principle make recurrence a question about electrical networks, and Lyons and
Peres's book \emph{Probability on Trees and Networks} (2016) develops the theory. The same cutsets in $\Z^3$ have about $24k^2$ edges, so the
lower bound converges, and indeed one can build a unit flow of finite energy to infinity: the walk on $\Z^3$ is transient, as Pólya
proved (Volume~3, Theorem~7.13). The monotonicity is not obvious from the walk: removing edges changes all the transition probabilities.
\end{solution}

\begin{solution}{pb:aldousbroder}
\idea Follow the tree of last exits along the walk: it is a Markov chain on rooted trees whose stationary law weighs each tree by the
product of its transition probabilities. Reversing time turns last exits into first entrances, and the root's weight cancels.

\writeup (a) The first visit to each $v\ne r$ comes after the first visit to some neighbour $u$, which was visited earlier, so the kept edge joins $v$ to an
earlier vertex. Each $v\ne r$ has one kept edge to an earlier vertex: there are $n-1$ edges and following them always decreases the time of first
visit, so they connect every vertex to $r$ without cycles. They form a spanning tree.

(b) In $T$, rooted at $x$, the vertex $y$ has one outgoing arc, and its subtree (the vertices whose path to $x$ passes through $y$) does not contain
$x$. Deleting $y$'s arc separates this subtree, now rooted at $y$, from the rest, rooted at $x$; adding $x\to y$ attaches the rest below $y$. The result is
an arborescence rooted at $y$. Each step sets the arc of the vertex being left to the vertex entered next, and later steps change it only
when that vertex is left again. So after a path that visits every vertex, the arc of every vertex other than the endpoint points to the
vertex that followed its last visit, whatever the starting state.

(c) Fix $(T',y)$. A predecessor $(T,x)$ must have $x\to y$ in $T'$, so $x$ is a child of $y$ in $T'$, and $T=T'-(x\to y)+(y\to z)$ for some $z$, which must lie in the
subtree of $x$ in $T'$ (so that $T$ is an arborescence rooted at $x$); conversely each such choice of $(x,z)$ with $z\sim y$ gives a predecessor. Then
$w(T,x)=w(T',y)\frac{\deg x}{\deg y}$ (the arc of $y$ is new, the arc of $x$ is gone), and the transition has probability $\frac1{\deg x}$. So
\[
\sum_{(T,x)\to(T',y)}w(T,x)\frac1{\deg x}=w(T',y)\sum_{x}\sum_{z}\frac1{\deg y}=w(T',y),
\]
since every neighbour $z$ of $y$ lies in the subtree of exactly one child $x$ of $y$. So $w$, hence $\mu$, is stationary. Next $w(T,x)=\frac{\deg x}{\prod_u\deg u}$, so
$Z=\sum_x\sum_Tw(T,x)=\frac{2|\mathcal E|\tau(G)}{\prod_u\deg u}$ (each spanning tree has one orientation towards each root) and $\mu(T,x)=\frac{\deg x}{2|\mathcal E|\tau(G)}$. Summing over $T$, the root has law
$\frac{\deg x}{2|\mathcal E|}=\pi(x)$.

(d) Let $(T_0,X_0)$ have law $\mu$ and run the tree chain along the walk $X_0,\dots,X_N$, which is the simple random walk started from $\pi$. Then $(T_N,X_N)$ has
law $\mu$. Let $A_N$ be the event that the walk visits every vertex by time $N$; on $A_N$, $T_N$ is the last-exit tree of $X_0,\dots,X_N$, rooted at $X_N$, by (b).
The reversed path $Y_s=X_{N-s}$ is again a simple random walk started from $\pi$ (Volume~3, Proposition~8.14), and the last visit of $X$ to $v$ before $N$
is the first visit of $Y$ to $v$; the vertex following it in $X$ precedes it in $Y$. So, as undirected trees, the last-exit tree of $X$ is the
first-entrance tree $F$ of $Y$, and $X_N=Y_0$. Hence for every tree $T$ and vertex $r$,
\[
\bigl|P(F=T,\ Y_0=r,\ A_N)-\mu(T,r)\bigr|\le P(A_N^c).
\]
As $N\to\infty$, $P(A_N^c)\to0$ (the cover time is finite almost surely, Problem~\ref{pb:commute}(d)), and on $A_N$ the first-entrance tree of $Y_0,\dots,Y_N$ is that of the
whole walk $Y$. So $P(F=T,\,Y_0=r)=\mu(T,r)$ and $P(F=T\mid Y_0=r)=\frac{\mu(T,r)}{\pi(r)}=\frac1{\tau(G)}$.

\tools{Volume~3, Proposition~8.14 (reversibility) and Corollary~8.15; Problem~\ref{pb:commute}(d).}
\leads The algorithm was found independently by Broder and by Aldous around 1989--1990; part (c) is the Markov chain tree theorem, which
holds for every finite irreducible chain with weights $\prod P(u,\text{parent}(u))$, and is a probabilistic twin of the matrix-tree theorem of Volume~11, Appendix~A.
Wilson's algorithm (1996), with loop-erased walks, samples a uniform spanning tree in the expected time of a few hitting times instead of
a cover time.
\end{solution}

\begin{solution}{pb:terras}
\idea The parity of $T^k(m)$ is a coin flip that has not yet been decided by the first $k$ bits of $m$: each new bit of $m$ decides exactly
one new parity.

\writeup (a) We prove by induction on $k$ that $v_k(m)$ depends only on $m\bmod2^k$ and that $T^k(m+2^k)=T^k(m)+3^j$, where $j$ is the number of ones in $v_k(m)$. For
$k=1$: $m$ and $m+2$ have the same parity, and $T(m+2)=T(m)+1$ if $m$ is even, $T(m)+3$ if $m$ is odd. Suppose it holds for $k$, and let $m'=m+2^{k+1}$.
Applying the statement twice, $v_k(m')=v_k(m)$ and $T^k(m')=T^k(m)+2\cdot3^j$, which has the same parity as $T^k(m)$; so $v_{k+1}(m')=v_{k+1}(m)$, and
$T^{k+1}(m')=T^{k+1}(m)+3^{j'}$ where $j'$ counts the ones in $v_{k+1}(m)$ (apply $T$ to two numbers of the same parity differing by $2\cdot3^j$). This proves the
first claim. For the bijection, it suffices to show injectivity on residues modulo $2^k$. If $m'=m+2^{k-1}$, the induction shows that
$v_{k-1}(m')=v_{k-1}(m)$ and $T^{k-1}(m')=T^{k-1}(m)+3^{j}$, an odd difference: the last parities differ. So the two residues modulo $2^k$ that lift a residue
modulo $2^{k-1}$ have different vectors, and by induction all $2^k$ vectors are different.

(b) A run of $2^k$ consecutive integers contains each residue modulo $2^k$ once, so $v_k(m)$ is uniform on $\{0,1\}^k$: its coordinates are independent
fair bits.

(c) Induction: if $T^i(m)=\lambda_im+\rho_i$, the next step multiplies by $\frac12$, or by $\frac32$ and adds $\frac12$, so $T^{i+1}(m)=\lambda_{i+1}m+\rho_{i+1}$ with $\rho_{i+1}\ge0$
depending only on the parities so far.

(d) The vector $(1,1,1,1,1)$ corresponds to one residue modulo $32$; since $T(m)=\frac{3m+1}2$ keeps $m\equiv-1\pmod{2^r}$ in the form $-1\pmod{2^{r-1}}$, it is
$m\equiv-1$, and $31,47,71,107,161$ are all odd: $m=31$. The vector of $1$ is $(1,0,1,0,1)$ (the cycle $1,2,1,2,\dots$), so the next integer with this vector is
$1+32=33$, whose trajectory $33,50,25,38,19$ indeed alternates.

\tools{Volume~1, Chapter~1 (congruences, induction).}
\leads Terras proved (a) in 1976, and Everett found it independently in 1977; Lagarias's survey of 1985 and his annotated bibliography
are the standard guides to the $3x+1$ problem. The bijection makes the random model of Problem~\ref{pb:everett} an exact statement about
the first $k$ steps.
\end{solution}

\begin{solution}{pb:everett}
\idea If at least $k(1-\frac{\ln2}{\ln3})$ of the first $k$ parities are even, the $k$ steps have multiplied by less than $1$. By Problem~\ref{pb:terras} the parities are fair
coins on residue classes, and the law of large numbers does the rest.

\writeup (a) By Problem~\ref{pb:terras}(c), for $m$ with $v_k(m)=u$, $T^k(m)=\lambda m+\rho_u$ with $\lambda=\frac{3^j}{2^k}<1$. Then $T^k(m)<m$ as soon as $m>\frac{\rho_u}{1-\lambda}$, so $\sigma(m)\le k$.

(b) Let $U_k$ be the set of $u$ with $3^j<2^k$, and $r_u$ the residue modulo $2^k$ with $v_k(r_u)=u$ (Problem~\ref{pb:terras}(a)). By (a) there is $m_0$ such that every $m>m_0$ with
$m\bmod2^k\in\{r_u:u\in U_k\}$ has $\sigma(m)\le k$. The density of that union of residue classes is $\frac{|U_k|}{2^k}=P(3^{B_k}<2^k)=P(B_k<k\frac{\ln2}{\ln3})$. So the lower density of
$\{\sigma\le k\}$ is at least this.

(c) Since $\frac{\ln2}{\ln3}=0.6309\ldots>\frac12$, the weak law of large numbers (Volume~3, Theorem~4.2) gives $P(B_k<k\frac{\ln2}{\ln3})\to1$. The set $\{\sigma<\infty\}$ contains $\{\sigma\le k\}$
for every $k$, so its lower density is $1$, and its density exists and equals $1$.

(d) In the model, $\ln T^{i+1}(m)-\ln T^i(m)$ is $\ln\frac12$ or approximately $\ln\frac32$ (the $+\frac12$ is negligible for large values), with probability $\frac12$ each:
mean $\frac12\ln\frac34\approx-0.144$. A typical $m$ therefore needs about $\frac{\ln m}{\frac12\ln(4/3)}\approx6.95\ln m$ steps of $T$ to come down to $1$. For $m$ near $10^6$ this
predicts $96$ steps; the actual average over $10^6\le m<10^6+2000$ is $92.8$. The prediction is a heuristic: the parities of later steps are
fair coins only in the sense of (a), for $k$ smaller than the number of digits of $m$.

\tools{Problem~\ref{pb:terras}; Volume~3, Theorem~4.2 (weak law of large numbers); natural density (Volume~1).}
\leads Tao proved in 2019 that almost all orbits (in logarithmic density) attain almost bounded values: for any function $f(m)\to\infty$, almost all $m$
have an iterate below $f(m)$. Whether every orbit reaches $1$ remains open; Volume~14, Appendix~F tells the story.
\end{solution}

\begin{solution}{pb:coprimewalk}
\idea Visibility is a sieve over primes. For each prime the walk equidistributes modulo $p$, except that $2$ is special: the parity of $X_t+Y_t$
is that of $t$. A uniform bound on the probability of divisibility by a large prime makes the infinite sieve converge.

\writeup (a) The steps $(\pm1,0)$ and $(0,\pm1)$ change $(U,V)$ by $(1,1)$, $(-1,-1)$, $(1,-1)$, $(-1,1)$, each with probability $\frac14$: independent fair signs.

(b) For odd $p$, $p\mid X$ and $p\mid Y$ if and only if $p\mid U$ and $p\mid V$, because $X=\frac{U+V}2$, $Y=\frac{U-V}2$ and $2$ is invertible modulo $p$. By independence the
probability is $P(p\mid U_t)^2$. The walk $U_t\bmod p$ on $\Z/p\Z$ with steps $\pm1$ is aperiodic for odd $p$ ($\{\pm1\}$ lies in no coset of the trivial subgroup, the only
proper one), so $P(p\mid U_t)\to\frac1p$ by Problem~\ref{pb:period}. For $p=2$: $X_t+Y_t=U_t\equiv t\pmod2$, so for odd $t$ one coordinate is odd. For even $t$, $U$ and $V$ are
even, and $X,Y$ are both even iff $U\equiv V\pmod4$. The walk $U_t\bmod4$ at even times moves on $\{0,2\}$ by steps $0,\pm2$ with probabilities
$\frac12,\frac14,\frac14$, so it converges to the uniform law on $\{0,2\}$, and so does $V$, independently: $P(U\equiv V\bmod4)\to\frac12$.

(c) $P(p\mid U_t)=\frac1p\sum_{k=0}^{p-1}E[e^{2\pi ikU_t/p}]=\frac1p\sum_{k=0}^{p-1}\cos^t\frac{2\pi k}p$. For $|y|\le\frac\pi2$, $\cos y=1-2\sin^2\frac y2\le1-\frac{2y^2}{\pi^2}\le e^{-2y^2/\pi^2}$, using $\sin\frac y2\ge\frac y\pi$; as $|\cos|$ has period
$\pi$, $|\cos x|\le e^{-2\delta(x)^2/\pi^2}$ where $\delta(x)$ is the distance from $x$ to $\pi\Z$. As $k$ runs over $\{1,\dots,p-1\}$, $\delta(\frac{2\pi k}p)$ takes each value $\frac{\pi j}p$, $1\le j\le\frac{p-1}2$, exactly
twice. So
\[
\sum_{k=0}^{p-1}\Bigl|\cos\frac{2\pi k}p\Bigr|^t\le1+2\sum_{j\ge1}e^{-2tj^2/p^2}\le1+2\int_0^\infty e^{-2tx^2/p^2}\dd x=1+p\sqrt{\frac\pi{2t}},
\]
and $P(p\mid U_t)\le\frac1p+\sqrt{\frac\pi{2t}}\le\frac1p+\frac2{\sqrt t}$.

(d) Let $E_p$ be the event that $p$ divides both coordinates; $\gcd(X_t,Y_t)=1$ means that no $E_p$ occurs and $(X_t,Y_t)\ne(0,0)$. Fix $P$ and let $Q=4\prod_{2<p\le P}p$. The pair
$(U_t\bmod Q,V_t\bmod Q)$ has independent coordinates, each of which is the $\pm1$ walk on $\Z/Q\Z$; this walk has period $2$ and at even (respectively odd) times converges to the uniform law on the
even (respectively odd) residues, by Problem~\ref{pb:period}. By the Chinese remainder theorem, along odd $t$ the events $E_p$, $p\le P$, become independent with
limits $P(E_2)=0$ and $P(E_p)=\frac1{p^2}$; along even $t$, $P(E_2)\to\frac12$ and $P(E_p)\to\frac1{p^2}$, again independently. So
\[
P\Bigl(\bigcap_{p\le P}E_p^c\Bigr)\to\alpha\prod_{2<p\le P}\Bigl(1-\frac1{p^2}\Bigr),\qquad\alpha=1\ (t\text{ odd}),\quad\alpha=\frac12\ (t\text{ even}).
\]
For the large primes: if $p>t$, $E_p$ forces $X_t=Y_t=0$, since $|X_t|,|Y_t|\le t$. For $P<p\le t$, by (c), $P(E_p)\le(\frac1p+\frac2{\sqrt t})^2\le\frac2{p^2}+\frac8t$. Hence
\[
0\le P\Bigl(\bigcap_{p\le P}E_p^c\Bigr)-P(\gcd(X_t,Y_t)=1)\le\sum_{p>P}\frac2{p^2}+\frac{8\pi(t)}t+P(X_t=Y_t=0),
\]
where $\pi(t)$ counts primes. As $t\to\infty$ the last two terms tend to $0$ (Chebyshev's bound gives $\pi(t)=O(\frac t{\ln t})$, Volume~1, Chapter~2), and then the first tends to $0$ as $P\to\infty$. Since $\prod_{p}(1-p^{-2})=\frac1{\zeta(2)}=\frac6{\pi^2}$ and the factor at $p=2$ is $\frac34$, $\prod_{p>2}(1-p^{-2})=\frac8{\pi^2}$. The limits are
$\frac8{\pi^2}\approx0.811$ for odd $t$ and $\frac4{\pi^2}\approx0.405$ for even $t$; one finds numerically $0.8108$ and $0.4055$ at $t=3001$ and $t=3000$.

\tools{Problem~\ref{pb:period}; Volume~1, Chapters~1--2 (Chinese remainder theorem, $\pi(t)=o(t)$, the Euler product for $\zeta(2)$); Volume~0 (Gaussian integral).}
\leads The average of the two limits is $\frac6{\pi^2}$, the density of visible points: the walk sees the plane as a uniform observer does, except for a
parity that it cannot forget. The same sieve works for any walk whose steps generate $\Z^2$, and the general phenomenon, an arithmetic
property of a random walk governed by the walk's behaviour modulo every integer, is the starting point of the theory of \emph{affine sieves}
of Bourgain, Gamburd and Sarnak.
\end{solution}

\begin{solution}{pb:range}
\idea A new point is discovered at time $s$ exactly when the walk, read backwards from time $s$, never comes back to its starting point.
So the expected range is a sum of non-return probabilities, which Volume~3 computed.

\writeup (a) $R_t=1+\sum_{s=1}^t\mathbf 1_{N_s}$, where $N_s$ is the event that $S_s\ne S_r$ for all $r<s$. Let $S'_q=S_s-S_{s-q}$ for $0\le q\le s$: it is again a simple walk (the steps in
reverse order), and $N_s$ says $S'_q\ne0$ for $1\le q\le s$. So $P(N_s)=P(S_1\ne0,\dots,S_s\ne0)$, and the term $s=0$ accounts for the $1$.

(b) Returns to $0$ only happen at even times, so the probability of no zero among times $1,\dots,2k+1$ equals that among times $1,\dots,2k$, which is
$u_{2k}$ by Volume~3, Theorem~7.3(a). Hence the terms $s=2k$ and $s=2k+1$ of (a) are both $u_{2k}$, and
\[
\mathbb E[R_{2m}]=2\sum_{k=0}^{m-1}u_{2k}+u_{2m}.
\]
Since $u_{2k+2}=u_{2k}\frac{2k+1}{2k+2}$, we have $(2m+1)u_{2m}=(2m+2)u_{2m+2}$, and the identity $\sum_{k=0}^mu_{2k}=(2m+1)u_{2m}$ follows by induction:
$(2m+1)u_{2m}+u_{2m+2}=(2m+3)u_{2m+2}$. So $\sum_{k=0}^{m-1}u_{2k}=2mu_{2m}$ and $\mathbb E[R_{2m}]=(4m+1)u_{2m}$. (Check: $m=1$ gives $\frac52$, the average of $2,3,3,2$.)

(c) $u_{2m}\sim\frac1{\sqrt{\pi m}}$ by Stirling's formula (Volume~0, Theorem~5.6), so
\[
\mathbb E[R_{2m}]\sim\frac{4m}{\sqrt{\pi m}}=\sqrt{\frac{16m}\pi}=\sqrt{\frac{8\cdot2m}\pi};
\]
odd times are squeezed between even ones since $R_t$ is non-decreasing.

(d) The same reversal gives
\[
\mathbb E[R_t]=\sum_{s=0}^tq_s,\qquad q_s=P(S_r\ne0\text{ for }1\le r\le s),
\]
and $q_s$ decreases to the probability of never returning. A Cesàro mean of a convergent sequence has the same limit.

\tools{Volume~3, Theorem~7.3 and Theorem~7.12; Stirling's formula (Volume~0, Theorem~5.6).}
\leads Dvoretzky and Erdős (1951) proved the law of large numbers $\frac{R_t}t\to1-F$ almost surely in $\Z^d$, $d\ge3$, and found the order $\frac t{\ln t}$ in the plane;
in $\Z^3$ the walk returns with probability about $0.34$, so it discovers about $0.66$ new points per step. The range is the simplest of the
\emph{occupation} questions that make random walks a laboratory for statistical physics.
\end{solution}

\begin{solution}{pb:cyclecover}
\idea The visited set grows as an interval, and each extension is a gambler's ruin. For the last vertex, only the two neighbours of the
candidate matter.

\writeup (a) The walker moves to adjacent vertices, so the visited set is connected: an arc. When it has $k<n$ vertices, the walker is at one of its
ends (it has just added it) or, at time $0$, at the unique vertex. Label the arc $1,\dots,k$ with the walker at $1$, and the two unvisited
neighbouring vertices $0$ and $k+1$ (distinct since $k<n-1$; if $k=n-1$ they coincide, which changes nothing below). A new vertex is added when the
walker leaves $\{1,\dots,k\}$, which takes $1\cdot(k+1-1)=k$ steps on average (Volume~3, Theorem~7.9).

(b) $\sum_{k=1}^{n-1}k=\frac{n(n-1)}2$.

(c) Fix $v\ne0$. Before it visits $v$, the walker must visit $v-1$ or $v+1$ (it starts at $0\ne v$, and $0$ may be one of them); let $T$ be the first time it is at
one of them, say at $v-1$ (the other case is symmetric). At time $T$, $v$ is unvisited. Now $v$ is the last vertex visited if and only if, from $v-1$,
the walker reaches $v+1$ before $v$, going the long way round. By the gambler's ruin on the arc from $v$ to $v$ through $v-1,\dots,v+1$, of length $n$, this has
probability $\frac1{n-1}$, by the strong Markov property at $T$.

(d) On $K_n$ the walk from $0$ is invariant under permutations of the other vertices, so the last one is uniform. On the star with centre $c$ and leaves
$\ell_1$, $\ell_2$, $\ell_3$, started at $\ell_1$, the centre is visited at time $1$, so the last vertex is $\ell_2$ or $\ell_3$, never $c$.

\tools{Volume~3, Theorem~7.9 (gambler's ruin) and the strong Markov property (Chapter~8).}
\leads Lovász and Winkler showed in 1993 that the cycle and the complete graph are the only graphs on which the last vertex is uniform from every
start. Cover times of large graphs are a deep subject: on the $n\times n$ torus the cover time is about $\frac4\pi n^2(\ln n)^2$ (Dembo, Peres, Rosen and
Zeitouni, 2004).
\end{solution}

\begin{solution}{pb:threegamblers}
\idea Polynomials in the fortunes whose expected change is constant turn expected durations into algebra. The product $XYZ$ loses on
average one third of the total per round, and the sum of pairwise products loses exactly one per round.

\writeup (a) Suppose all three are in the game. If the pair $(X,Y)$ plays, the change is $(X+\varepsilon)(Y-\varepsilon)Z-XYZ=(\varepsilon(Y-X)-\varepsilon^2)Z$ with $\varepsilon=\pm1$ equally likely,
of conditional mean $-Z$. Averaging over the three pairs, $\mathbb E[\Delta(XYZ)\mid\mathcal F_t]=-\frac{X+Y+Z}3=-\frac s3$. So $M_t=X_tY_tZ_t+\frac s3t$, stopped at the first ruin
$\tau$, is a martingale. The state space is finite and from every state a ruin can be reached within $s$ rounds, so $\mathbb E[\tau]<\infty$ (Volume~3, Lemma~7.8), and the
increments are bounded; optional stopping (Volume~3, Theorem~10.12) gives $abc=\mathbb E[M_\tau]=0+\frac s3\mathbb E[\tau]$.

(b) With three players, the pair $(X,Y)$ changes $XY+YZ+ZX$ by $\varepsilon(Y-X)-\varepsilon^2+\varepsilon Z-\varepsilon Z$, of mean $-1$, and similarly for the other pairs. With two players,
say $Z=0$, only $XY$ changes, by $\varepsilon(Y-X)-1$, of mean $-1$. So $Q_t+t$ is a martingale until the end $\sigma$, where $Q=XY+YZ+ZX$ is $0$. As in (a), $\mathbb E[\sigma]<\infty$ and
$\mathbb E[\sigma]=Q_0=ab+bc+ca$.

(c) $X_t$ is a bounded martingale (each round is fair for each player, whether or not the player plays), so $a=\mathbb E[X_\sigma]=s\,P(\text{first player wins all})$.

(d) With $a=b=c=4$: the game lasts $48$ rounds on average, and the first player leaves after $16$ rounds on average.

\tools{Volume~3, Theorem~10.12 (optional stopping) and Lemma~7.8.}
\leads For two players the same martingale $XY+t$ is the gambler's ruin duration $ab$ of Volume~3, Theorem~7.9. The probability that a given player is
the first to be ruined is a much harder quantity: it is a harmonic measure for a random walk in a triangle. The analogue with $n$ players is the expected time to absorption of a random walk in a simplex.
\end{solution}

\begin{solution}{pb:moran}
\idea Each step is fair, so the type frequency is a martingale; the only randomness is in the speed, which is slowest when one type is
rare. Counting moves state by state reduces the time to a Green function of the simple walk.

\writeup (a) From $x$, the walk goes up if an $A$ reproduces and a $B$ dies, probability $\frac{x(N-x)}{N^2}$, and down with the same probability. So $X$ is a bounded
martingale, absorbed at $0$ or $N$ in finite expected time (Volume~3, Lemma~7.8), and $i=\mathbb E[X_\tau]=N\,P(X_\tau=N)$.

(b) With $\delta=X_{t+1}-X_t$: $X_{t+1}(N-X_{t+1})=X_t(N-X_t)+\delta(N-2X_t)-\delta^2$, and $\mathbb E[\delta\mid X_t]=0$, $\mathbb E[\delta^2\mid X_t]=\frac{2X_t(N-X_t)}{N^2}$.

(c) Let $q_x=\frac{2x(N-x)}{N^2}$, the probability of moving from $x$. The sequence of positions after each move is the simple random walk on $\{0,\dots,N\}$ absorbed at the ends,
and each stay at $x$ before a move lasts on average $\frac1{q_x}$ steps (geometric). Let $G(i,x)$ be the expected number of moves from $x$ made by the simple walk
started at $i$, that is its expected number of visits to $x$. For $x\ge i$ the walk reaches $x$ with probability $\frac ix$ (gambler's ruin). From $x$ it escapes to $\{0,N\}$
before returning to $x$ with probability $\frac12\bigl(\frac1x+\frac1{N-x}\bigr)=\frac N{2x(N-x)}$, so the number of visits, once $x$ is reached, is geometric with mean $\frac{2x(N-x)}N$. Multiplying, $G(i,x)=\frac ix\cdot\frac{2x(N-x)}N=\frac{2i(N-x)}N$ for $x\ge i$, and symmetrically $G(i,x)=\frac{2x(N-i)}N$ for $x\le i$. Hence
\[
\begin{aligned}t_i=\sum_{x=1}^{N-1}\frac{G(i,x)}{q_x}&=\sum_{x\ge i}\frac{2i(N-x)}N\cdot\frac{N^2}{2x(N-x)}+\sum_{x<i}\frac{2x(N-i)}N\cdot\frac{N^2}{2x(N-x)}\\&=Ni\sum_{x=i}^{N-1}\frac1x+N(N-i)\sum_{x=1}^{i-1}\frac1{N-x},\end{aligned}
\]
and the last sum is $\sum_{y=N-i+1}^{N-1}\frac1y$.

(d) For $i=1$ the second sum is empty and $t_1=NH_{N-1}\sim N\ln N$. For $i=\lfloor pN\rfloor$, $\sum_{x=i}^{N-1}\frac1x=\ln\frac Ni+O(\frac1i)=\ln\frac1p+o(1)$ and $\sum_{y=N-i+1}^{N-1}\frac1y=\ln\frac1{1-p}+o(1)$, so
$t_i=N^2\bigl(p\ln\frac1p+(1-p)\ln\frac1{1-p}+o(1)\bigr)$.

\tools{Volume~3, Theorems~7.9 and 10.12, Lemma~7.8; harmonic sums (Volume~0).}
\leads Moran introduced the model in 1958 as a tractable cousin of the Wright--Fisher model, in which the whole generation is resampled at
once; the diffusion approximation of Wright and Kimura gives the same entropy-like formula for that model, with time counted in
generations. The decay of heterozygosity in (b) is the population geneticist's measure of genetic drift. Kingman's coalescent (1982) runs the same
models backwards in time, following the ancestry of a sample.
\end{solution}

\begin{solution}{pb:bernoullilaplace}
\idea The chain maps polynomials of degree $k$ to polynomials of degree $k$, so its matrix is triangular in the basis of powers, and the
eigenvalues are read on the diagonal.

\writeup (a) From $x$: the number of red balls in the first urn drops if a red is drawn there ($\frac xn$) and a white in the second (the second urn
has $n-x$ red and $x$ white): $P(x,x-1)=\frac{x^2}{n^2}$. It rises with probability $P(x,x+1)=\frac{(n-x)^2}{n^2}$, and stays otherwise. Detailed balance:
$\pi(x)P(x,x+1)=\binom nx^2\frac{(n-x)^2}{n^2}\binom{2n}n^{-1}$ and $\pi(x+1)P(x+1,x)=\binom n{x+1}^2\frac{(x+1)^2}{n^2}\binom{2n}n^{-1}$, equal since $\binom nx(n-x)=\binom n{x+1}(x+1)$; and $\sum_x\binom nx^2=\binom{2n}n$.

(b) With $\delta=X_1-x\in\{-1,0,1\}$: $\mathbb E[\delta]=\frac{(n-x)^2-x^2}{n^2}=\frac{n-2x}n$ and $\mathbb E[\delta^2]=\frac{x^2+(n-x)^2}{n^2}$. Then $\mathbb E[(x+\delta)^k]=\sum_r\binom krx^{k-r}\mathbb E[\delta^r]$, and $\mathbb E[\delta^r]$ is
$\mathbb E[\delta]$ or $\mathbb E[\delta^2]$, a polynomial in $x$ of degree $1$ or $2$. So the terms $r\ge3$ have degree at most $k-1$, and the coefficient of $x^k$ in $P(x^k)$ is
$1+k\cdot\bigl(-\frac2n\bigr)+\binom k2\frac2{n^2}=1-\frac{2k}n+\frac{k(k-1)}{n^2}=1-\frac{k(2n-k+1)}{n^2}$. Functions on $\{0,\dots,n\}$ are exactly the polynomials of degree at most $n$ (Lagrange interpolation),
and in the basis $1,x,\dots,x^n$ the operator $P$ is triangular with diagonal $\lambda_0,\dots,\lambda_n$. These are its eigenvalues.

(c) $\mathbb E[X_{t+1}-\frac n2\mid X_t]=X_t-\frac n2+\frac{n-2X_t}n=(1-\frac2n)(X_t-\frac n2)$, so $\mathbb E[X_t]-\frac n2=\frac n2(1-\frac2n)^t$. For $0\le y\le\frac12$, $\ln(1-y)\ge-y-y^2$, so for $n\ge4$ and $t\le\frac n4\ln n$,
$t\ln(1-\frac2n)\ge-\frac n4\ln n\,(\frac2n+\frac4{n^2})=-\frac12\ln n-\frac{\ln n}n$, and $\frac n2(1-\frac2n)^t\ge\frac{\sqrt n}2e^{-\ln n/n}\ge\frac{\sqrt n}2e^{-1/e}>\frac{\sqrt n}3$. Under $\pi$, $X$ is hypergeometric (draw $n$
balls from $2n$, half red), with variance $\frac{n^2}{4(2n-1)}<\frac n4$.

\tools{Volume~3, Proposition~8.14 (detailed balance); polynomial interpolation (Volume~0).}
\leads Daniel Bernoulli posed the model in 1769 and Laplace solved its diffusion limit in his \emph{Théorie analytique des probabilités} (1812). Diaconis
and Shahshahani (1987) used the eigenvalues above (with their eigenvectors, the dual Hahn polynomials) to show that the chain mixes after
$\frac n4\ln n$ steps and not before, a cutoff; (c) is the lower half of their theorem.
\end{solution}

\begin{solution}{pb:binompowers}
\idea Raising the binomial law to the power $s$ keeps it Gaussian, with variance divided by $s$. Near the middle the local limit theorem applies; away from it the terms
are too small to matter.

\writeup Let $p_k=\binom nk2^{-n}$, $x_k=\frac{2k-n}{\sqrt n}$ and $\phi(x)=\frac{e^{-x^2/2}}{\sqrt{2\pi}}$, so the sum is $2^{sn}\Sigma$ with $\Sigma=\sum_kp_k^s$. Fix $A>0$. By Volume~3, Theorem~6.12 (with $p=\frac12$,
$\sqrt{npq}=\frac{\sqrt n}2$), $p_k=\frac2{\sqrt n}\phi(x_k)(1+\varepsilon_k)$ with $\max_{|x_k|\le A}|\varepsilon_k|\to0$. The points $x_k$ are spaced $\frac2{\sqrt n}$ apart, so by Riemann sums
\[
\Bigl(\frac{\sqrt n}2\Bigr)^{s-1}\sum_{|x_k|\le A}p_k^s=\bigl(1+o(1)\bigr)\frac2{\sqrt n}\sum_{|x_k|\le A}\phi(x_k)^s\to\int_{-A}^A\phi^s.
\]
For the rest, $p_k\le p_{\lfloor n/2\rfloor}\le\frac c{\sqrt n}$ (central binomial coefficient), and $\sum_{|x_k|>A}p_k\le\frac1{A^2}$ by Chebyshev, since $x_K$ has mean $0$ and variance $1$ when $K$ is binomial.
So $(\frac{\sqrt n}2)^{s-1}\sum_{|x_k|>A}p_k^s\le(\frac c2)^{s-1}A^{-2}$. Letting $n\to\infty$, then $A\to\infty$,
\[
\Bigl(\frac{\sqrt n}2\Bigr)^{s-1}\Sigma\to\int_\R\phi^s=(2\pi)^{-s/2}\sqrt{\frac{2\pi}s}=\frac{(2\pi)^{(1-s)/2}}{\sqrt s}.
\]
Hence $\Sigma\sim(\frac2{\sqrt n})^{s-1}\frac{(2\pi)^{(1-s)/2}}{\sqrt s}=\frac1{\sqrt s}(\frac2{\pi n})^{(s-1)/2}$. For $s=2$ this gives $\frac{4^n}{\sqrt2}\sqrt{\frac2{\pi n}}=\frac{4^n}{\sqrt{\pi n}}$, which is $\binom{2n}n$ asymptotically (by
Vandermonde, $\sum\binom nk^2=\binom{2n}n$). For $s=3$: $\sum\binom nk^3\sim\frac{2\cdot8^n}{\sqrt3\,\pi n}$; for $s=4$: $\sum\binom nk^4\sim\frac{16^n}2\bigl(\frac2{\pi n}\bigr)^{3/2}$. Numerically, one finds ratios $0.9997$
and $0.9994$ at $n=1000$.

\tools{Volume~3, Theorem~6.12 (local de Moivre--Laplace) and Chapter~4 (Chebyshev); Volume~4, Chapter~4 (Gaussian integral).}
\leads For $s=3$ these are the Franel numbers; Franel (1894, 1895) showed that the sums for $s=3$ and $s=4$ satisfy linear recurrences with polynomial coefficients, and this
asymptotic is what such a recurrence must reproduce. The same local-limit argument gives $\sum_kP(S_n=k)^s$ for any lattice random walk.
\end{solution}

\begin{solution}{pb:saw}
\idea A long self-avoiding walk is two shorter ones glued together, but not every gluing is self-avoiding: submultiplicativity, hence a growth rate by Fekete. Simple
families of walks squeeze it.

\writeup (a) Cutting a self-avoiding walk of length $m+n$ after $m$ steps gives a self-avoiding walk of length $m$ and, translated to the origin, one of length $n$; the map is injective.

(b) Fix $k$ and write $n=qk+r$ with $0\le r<k$. Iterating, $a_n\le qa_k+a_r$ (with $a_0$ replaced by $0$ when $r=0$). So $\frac{a_n}n\le\frac{qk}n\cdot\frac{a_k}k+\frac{\max_{r<k}|a_r|}n$, and $\limsup\frac{a_n}n\le\frac{a_k}k$ for every
$k$. Hence $\limsup\frac{a_n}n\le\inf_k\frac{a_k}k\le\liminf\frac{a_n}n$.

(c) Apply (b) to $a_n=\ln c_n$.

(d) After the first step a self-avoiding walk never steps straight back, so $c_n\le4\cdot3^{n-1}$ and $\mu\le3$. For the lower bound, a walk with steps east, north, south and no north--south
reversal is self-avoiding: its $x$-coordinate never decreases, and at a given $x$ it moves monotonically along one vertical segment. Let $e_n$ and $v_n$ count such walks of length
$n$ ending with a horizontal or vertical step. Then $e_{n+1}=e_n+v_n$ and $v_{n+1}=2e_n+v_n$, with $e_1=1$, $v_1=2$. The matrix $\begin{psmallmatrix}1&1\\2&1\end{psmallmatrix}$ has eigenvalues $1\pm\sqrt2$, with eigenvectors $(1,\pm\sqrt2)$, and
$(1,2)=\frac{1+\sqrt2}2(1,\sqrt2)+\frac{1-\sqrt2}2(1,-\sqrt2)$ has a positive component along the first. So $e_n+v_n\ge c(1+\sqrt2)^n$
with $c>0$, and $\mu\ge1+\sqrt2$.

(e) $\mu\le c_{14}^{1/14}=2374444^{1/14}=2.854$.

\tools{Volume~0, Chapter~7 (eigenvalues of $2\times2$ matrices); Volume~0, Chapter~5 (lim sup and lim inf).}
\leads Hammersley and Morton (1954) proved that $\mu$ exists. Its value on $\Z^2$ is unknown; numerical work gives $\mu=2.6381585\ldots$ On the hexagonal lattice, Nienhuis (1982) predicted
$\mu=\sqrt{2+\sqrt2}$, and Duminil-Copin and Smirnov (2012) proved it. The ratios $\frac{c_{n+1}}{c_n}$ in Figure~\ref{fig:saw} oscillate with the parity of $n$ but approach $\mu$ much faster than $c_n^{1/n}$.
\end{solution}

\begin{solution}{pb:inspection}
\idea A random integer is more likely to fall in a long gap than in a short one, in proportion to the length. So the gap seen from a random point has mean $\frac{\sum g^2}{\sum g}$,
and Cauchy--Schwarz already makes this at least the average gap.

\writeup (a) Let $N=\pi(x)$. The gaps with $p_{n+1}\le x$ are $g_1,\dots,g_{N-1}$, with sum $p_N-2$. By Cauchy--Schwarz,
$\sum_{n<N}g_n^2\ge\frac{(p_N-2)^2}{N-1}$. By the prime number theorem $N\sim\frac x{\log x}$, and $p_N\sim x$ (since $p_{N+1}>x$ and $p_{N+1}\sim p_N$). Hence $\sum g_n^2\ge(1+o(1))x\log x$.

(b) Each gap $]p_n,p_{n+1}]$ contains $g_n$ integers, each of which sees $G=g_n$. So $\sum_{2<m\le p_N}G(m)=\sum_{n<N}g_n^2$, and dividing by $p_N-2$ gives $\mathbb E[G]$; by (a) it is at least
$(1+o(1))\frac{x\log x}x=(1+o(1))\log x$. The average gap $\frac{p_N-2}{N-1}$ is $\sim\log x$: the gap containing a random point is size-biased, as the bus interval containing a random
arrival time in Volume~3, Proposition~9.10.

(c) With $P(G_0=k)=q(1-q)^{k-1}$, $\mathbb E[G_0]=\frac1q$ and $\mathbb E[G_0(G_0-1)]=\frac{2(1-q)}{q^2}$, so $\mathbb E[G_0^2]=\frac{2-q}{q^2}$ and the ratio is $\frac{2-q}q$. With $q=\frac1{\log x}$,
the size-biased mean is $\sim2\log x$, and $\sum g_n^2\approx(\text{number of gaps})\cdot\mathbb E[G_0^2]\approx\frac x{\log x}\cdot2\log^2x=2x\log x$.

(d) The ratio climbs slowly, $1.16$, $1.52$, $1.63$, well above the proved bound $1$ and still well below $2$. It is consistent with the prediction $2$ approached with corrections of order
$\frac1{\log x}$ (the gaps near $10^8$ are not yet exponential: gaps of $2$ and $4$ are over-represented and the mean gap $\log x$ is only $18$). If one draws $m$ at random and takes the next prime,
the primes that follow long gaps are favoured, with probability proportional to the preceding gap: about twice the uniform weight on average, if Cram\'er is right.

\tools{Volume~3, Proposition~2.6 (Cauchy--Schwarz) and Proposition~9.10 (inspection paradox); Volume~5, Theorem~5.5 and Chapter~10.}
\leads Some prime generators used in cryptography draw a random number and search upwards for the next prime (Brandt and Damg\aa rd, 1992). By (b) their output is not uniform among
primes; Fouque and Tibouchi (2014) showed, for instance, that the upper member of a twin pair is output much less often than a uniform choice would. Proving
$\sum_{p_{n+1}\le x}g_n^2=O(x^{1+\varepsilon})$ without the Riemann hypothesis is an open problem.
\end{solution}

\section*{Sources and further reading}

\begin{itemize}
\item N.~Alon and J.~Spencer, \emph{The Probabilistic Method}, 4th ed., Wiley, 2016. The reference for the whole chapter.
\item M.~Aigner and G.~Ziegler, \emph{Proofs from THE BOOK}, 6th ed., Springer, 2018. Chapters on the probabilistic method and
on the crossing lemma.
\item P.~Erd\H os, \emph{Some remarks on the theory of graphs}, Bulletin of the American Mathematical Society 53 (1947),
292--294.
\item P.~Erd\H os and A.~R\'enyi, \emph{On the evolution of random graphs}, Publications of the Mathematical Institute of the
Hungarian Academy of Sciences 5 (1960), 17--61.
\item L.~Lov\'asz and P.~Erd\H os, \emph{Problems and results on 3-chromatic hypergraphs and some related questions}, in
Infinite and Finite Sets, Colloquia Mathematica Societatis J\'anos Bolyai 10, North-Holland, 1975, 609--627.
\item R.~Moser and G.~Tardos, \emph{A constructive proof of the general Lov\'asz local lemma}, Journal of the ACM 57 (2010),
article~11.
\item M.~Krivelevich and B.~Sudakov, \emph{The phase transition in random graphs: a simple proof}, Random Structures and
Algorithms 43 (2013), 131--138.
\item S.~Janson, T.~\L uczak and A.~Ruci\'nski, \emph{Random Graphs}, Wiley, 2000; A.~Frieze and M.~Karo\'nski, \emph{Introduction to
Random Graphs}, Cambridge University Press, 2016; R.~van der Hofstad, \emph{Random Graphs and Complex Networks}, Vol.~1,
Cambridge University Press, 2017.
\item M.~Campos, S.~Griffiths, R.~Morris and J.~Sahasrabudhe, \emph{An exponential improvement for diagonal Ramsey},
preprint, 2023, arXiv:2303.09521.
\item J.~H\aa stad, \emph{Some optimal inapproximability results}, Journal of the ACM 48 (2001), 798--859.
\item M.~Mitzenmacher and E.~Upfal, \emph{Probability and Computing}, 2nd ed., Cambridge University Press, 2017, Chapter~6.
\end{itemize}
